\documentclass[12pt,a4paper,UTF8]{ctexart}

\usepackage[a4paper,top=2.25cm,bottom=2.25cm,left=2.25cm,right=2.25cm]{geometry}
\usepackage{amsmath,amssymb,mathtools}
\usepackage{booktabs,longtable,tabularx,array,multirow}
\usepackage{graphicx}
\usepackage{xcolor}
\usepackage{listings}
\usepackage{enumitem}
\usepackage{float}
\usepackage{tikz}
\usepackage{hyperref}
\usepackage{url}
\usetikzlibrary{arrows.meta,positioning,shapes.geometric,fit,calc,backgrounds}

\definecolor{CampBlue}{RGB}{34,92,157}
\definecolor{CampGreen}{RGB}{55,132,94}
\definecolor{CampOrange}{RGB}{215,123,49}
\definecolor{CampPurple}{RGB}{119,83,157}
\definecolor{CampGray}{RGB}{245,247,250}
\definecolor{CampRed}{RGB}{179,53,63}

\hypersetup{
  colorlinks=true,
  linkcolor=CampBlue,
  citecolor=CampGreen,
  urlcolor=CampPurple,
  pdftitle={CAMP-VS 虚拟染色竞赛项目技术报告},
  pdfauthor={CAMP-VS Project}
}

\lstdefinestyle{shell}{
  basicstyle=\ttfamily\small,
  backgroundcolor=\color{CampGray},
  frame=single,
  rulecolor=\color{black!20},
  breaklines=true,
  columns=fullflexible,
  keepspaces=true,
  showstringspaces=false,
  aboveskip=0.8em,
  belowskip=0.8em
}
\lstdefinestyle{python}{
  language=Python,
  basicstyle=\ttfamily\small,
  keywordstyle=\color{CampBlue}\bfseries,
  commentstyle=\color{CampGreen},
  stringstyle=\color{CampOrange},
  backgroundcolor=\color{CampGray},
  frame=single,
  rulecolor=\color{black!20},
  breaklines=true,
  columns=fullflexible,
  showstringspaces=false
}
\lstdefinestyle{yaml}{
  basicstyle=\ttfamily\small,
  keywordstyle=\color{CampBlue}\bfseries,
  commentstyle=\color{CampGreen},
  stringstyle=\color{CampOrange},
  backgroundcolor=\color{CampGray},
  frame=single,
  rulecolor=\color{black!20},
  breaklines=true,
  columns=fullflexible,
  keepspaces=true,
  showstringspaces=false,
  morecomment=[l]{\#},
  morestring=[b]"
}

\newcommand{\code}[1]{\texttt{\detokenize{#1}}}
\newcommand{\important}[1]{%
  \begin{center}
  \fbox{\begin{minipage}{0.92\textwidth}\small #1\end{minipage}}
  \end{center}}
\newcommand{\term}[2]{\textbf{#1（#2）}}

\setlist[itemize]{itemsep=0.35em,topsep=0.4em}
\setlist[enumerate]{itemsep=0.35em,topsep=0.4em}
\setcounter{tocdepth}{2}
\setcounter{secnumdepth}{3}
\emergencystretch=2em

\title{\textbf{CAMP-VS 虚拟染色竞赛项目技术报告}\\[0.6em]
\large 从 DAPI 图像到 CD68 虚拟染色的完整工程解析}
\author{基于当前项目源码、配置与本地官方数据审计结果整理}
\date{源码核对日期：2026 年 8 月 24 日}

\begin{document}

\maketitle

\begin{abstract}
本文面向第一次接触医学图像、深度学习和命令行的读者，完整说明 CAMP-VS
虚拟染色竞赛工程。项目把一张 DAPI 荧光显微镜图像作为输入，预测同一位置的
CD68 标记图像。本文从生物学背景讲起，逐步解释数据目录、ROI 分组、manifest、
PyTorch DataLoader、Sobel 结构提示、四尺度 NAF 编解码器、余弦原型混合、任务适配器、
重型四尺度 Detail U-Net、自适应空间--频率路线、三路内容感知融合、真实 $3\times3$ ROI 邻域、两阶段复合损失、
混合精度训练、EMA、三域验证、D4 测试时增强和提交 ZIP 的产生过程。

报告不是根据模型名称猜测而写，而是逐项核对了当前工作区中的
\code{README.md}、\code{docs/essay/report.tex}、默认/初赛/v2/v3/v4/resolved-local YAML 配置、CLI、训练器、
验证器、模型、损失、数据集、指标及提交代码，并读取了本地官方数据审计结果。
因此，本文会特别指出三项容易被误解的事实：第一，本地 JPG 在磁盘上是 RGB，
只是三个通道逐像素完全相同；第二，
\code{configs/resolved_local.yaml} 已把这种审计结论解析为 1 个逻辑输入/输出通道，
但提交时仍恢复为原始 RGB 存储语义；第三，本地 JPG round-trip 指标是接近
最终文件形态的代理指标，但没有官方评测器源码时，不能宣称它与平台算法完全相同。
\end{abstract}

\important{\textbf{阅读方法。} 如果你只学过计算机专业的机器学习和深度学习课程，请不要跳过第 1 节。
这一节先建立“组织--ROI--patch--细胞--像素”的共同语言，再解释五个 marker、荧光强度、监督标签和
评价层级；理解这些概念后，后面的网络结构和损失函数才不只是公式。如果你只想在 AutoDL 上跑通项目，
可直接阅读附录“AutoDL 从上传到训练的保姆级流程”。文中“稳定对照”指
\code{configs/initial_round_cd68_retrain_v2.yaml}；“v3”是已停止的轻量多分支候选；“当前候选”或“v4”专指
\code{configs/initial_round_cd68_max_v4.yaml} 与默认/resolved-local 配置合并后的实际设置。}

\tableofcontents
\listoffigures
\listoftables
\clearpage

\section{引言与项目背景}

\subsection{什么是染色，为什么显微镜图像需要它}

人体组织中的细胞通常接近透明。直接把一块薄薄的组织放到显微镜下，很多结构并不会自动呈现出
清晰的颜色。实验人员会使用能够与特定分子结合的染料或抗体，把原本不明显的结构“点亮”。
这件事可以类比为给城市地图加图例：道路、河流、医院本来都在地图上，但用不同颜色标出后，
人才能快速分辨它们。

\term{免疫组织化学染色}{Immunohistochemistry, IHC}利用抗体识别组织中的目标蛋白。
传统 IHC 往往需要制备切片、加入抗体试剂、孵育、清洗和显色，过程可能持续数小时，试剂和人工
都有成本，而且一次实验会对样本产生不可逆处理。若希望观察多个标记物，通常还要进行更多实验。

\subsection{H\&E、IHC、免疫荧光不是同一种数据}

初学者容易把所有“有颜色的病理图”都叫染色图，但它们的成像机制和计算含义不同。表~\ref{tab:stain-types}
给出最小区分。本项目文件夹按 DAPI 与四个蛋白 marker 分通道保存，像素近似灰度强度；在缺少更完整的
官方实验协议时，报告把它们作为\textbf{配准的荧光强度通道}处理，不自行猜测具体显微镜型号、荧光团或
抗体克隆。

\begin{table}[H]
\centering
\caption{常见组织成像方式与本项目数据的区别}
\label{tab:stain-types}
\begin{tabularx}{\textwidth}{>{\raggedright\arraybackslash}p{0.14\textwidth}
                              >{\raggedright\arraybackslash}p{0.22\textwidth}
                              >{\raggedright\arraybackslash}p{0.25\textwidth}X}
\toprule
方式 & 实验上显示什么 & 计算机看到什么 & 与本项目的关系 \\
\midrule
H\&E & 苏木精偏向细胞核、伊红显示胞质和基质，是常规组织形态染色 & 通常是彩色 RGB 图，颜色由两种染料混合 & 本项目没有 H\&E 输入，不能假设网络看到了其丰富组织颜色。\\
IHC & 抗体识别目标蛋白，再用显色反应显示阳性区域 & 常见为明场彩色图，需处理染料颜色分离 & CD68 等也常用于 IHC，但本项目磁盘图是独立 marker 强度通道，不应套用普通彩色 IHC 的颜色规则。\\
免疫荧光 & 抗体携带/连接荧光信号，不同分子由不同通道记录 & 每个通道本质是空间强度矩阵，可单独灰度显示，也可伪彩叠加 & 本项目的数据组织最接近这种“DAPI 加多个 marker 通道”的计算形式。\\
虚拟染色 & 不新增化学探针，由模型预测目标通道 & 输入张量经过神经网络得到连续强度图 & 预测图只是对实验图的统计估计，不能反过来当作真实做过一次染色。\\
\bottomrule
\end{tabularx}
\end{table}

\subsection{虚拟染色是什么}

\term{虚拟染色}{Virtual Staining}并不是让计算机真的把染料滴到玻片上，而是训练一个人工智能模型：
给它一种已经拍摄好的图像，让它预测同一位置在另一种标记下应该呈现的像素强度。对本项目来说，
可以把模型想象成一个非常复杂的“逐像素翻译器”：输入一张 DAPI 图，输出一张 CD68 图，输入和输出
的尺寸、位置一一对应。

它的潜在价值包括：

\begin{itemize}
  \item 减少重复染色所需的时间和试剂；
  \item 在不再次处理样本的前提下，探索另一种标记物的可能分布；
  \item 同一张 DAPI 图理论上可作为多个目标 marker 的共同输入；
  \item 预测是软件过程，可以批量运行并保留可复现的参数与日志。
\end{itemize}

需要强调的是，竞赛模型输出是\textbf{算法预测}，不是经过临床验证的真实检测结果。它可以用于竞赛、
研究和方法开发，但不能因为图像“看起来像”就直接替代病理诊断。

\subsection{DAPI 与 CD68 分别是什么}

\term{DAPI}{4',6-diamidino-2-phenylindole，4',6-二脒基-2-苯基吲哚}是一种常用的蓝色荧光 DNA
染料。它与双链 DNA 的富含 A--T 区域结合后荧光明显增强，所以在荧光显微镜下常被用作细胞核复染。
简单地说，DAPI 图像主要告诉我们“细胞核在哪里、密度如何、形状怎样”。DAPI 的常见性质可参考
Thermo Fisher 的技术说明\cite{thermofisher_dapi}。

\term{CD68}{Cluster of Differentiation 68}是一种主要位于内体/溶酶体系统的糖蛋白，在单核细胞和
组织巨噬细胞中常见。病理研究中常用 CD68 抗体辅助观察巨噬细胞相关信号和炎症细胞浸润；Human
Protein Atlas 也给出了相应的组织表达资料\cite{hpa_cd68}。但是 CD68 并非“只在巨噬细胞出现”的
绝对身份证，一些非髓系细胞也可能表达，因此单凭 CD68 阳性不能完成细胞身份判定\cite{cd68_specificity}。
为了入门，可以暂时把 CD68 图理解成“巨噬细胞相关信号的空间热力图”，同时牢记这只是近似说法。

只看 DAPI 并不能直接读取 CD68。模型需要从大量同位置的 DAPI--CD68 配对样本中学习统计规律：
哪些核形态、密度和周围纹理经常与 CD68 信号共同出现。这个映射不是确定的化学公式，因此虚拟染色
本质上是一个有不确定性的预测问题。

\subsection{荧光显微镜图像与 patch}

\term{荧光显微镜图像}{Fluorescence Microscopy Image}记录荧光分子在不同位置发出的光强。
存储到文件后，每个像素是一个数字：数值较大通常代表该位置的信号较强。项目内部会把 8 位像素
$0\sim255$ 转成浮点数 $0\sim1$，以便神经网络计算。

一张完整组织切片可能非常大，不能总是整张送入 GPU。因此数据制作者会把大图裁成许多固定大小的
\term{patch}{图块}。本地官方数据的每个 patch 是 $256\times256$ 像素。文件名如
\code{ROI000_03_12.jpg}，可拆成：

\begin{itemize}
  \item \code{ROI000}：第 0 个感兴趣区域；
  \item \code{03}：该 patch 在大区域网格中的行坐标；
  \item \code{12}：列坐标；
  \item \code{.jpg}：JPEG 文件格式。
\end{itemize}

\term{ROI}{Region of Interest，感兴趣区域}可以理解为从大切片中选出的一块连续区域。同一 ROI
里的相邻 patch 可能原本就在大图上彼此相邻，因此它们并不是完全独立的样本。

\subsection{竞赛目标与图像到图像翻译}

项目的核心任务是：对每张测试 DAPI patch 生成同尺寸的 CD68 patch。模型学习的是

\begin{equation}
  f_{\theta}:\; \mathbf{x}_{\mathrm{DAPI}}\longrightarrow
  \widehat{\mathbf{y}}_{\mathrm{CD68}},
\end{equation}

其中 $\theta$ 是模型参数，$\mathbf{x}$ 是输入图，$\widehat{\mathbf{y}}$ 是预测图。因为输入和输出都是
图像，而且像素位置需要对应，所以它属于\term{图像到图像翻译}{Image-to-Image Translation}，
更具体地说也属于图像恢复/回归问题。

本地验证主要观察 \term{SSIM}{Structural Similarity，结构相似性}与
\term{PSNR}{Peak Signal-to-Noise Ratio，峰值信噪比}。SSIM 更关注亮度、对比度和局部结构是否相似，
PSNR 则由均方误差推导。两者通常都是越高越好，但它们不能替代医学专家判断。

\subsection{先建立空间层级：组织、ROI、patch、细胞和像素}

计算机视觉教材常从一张独立的图像讲起；医学显微图像却有天然的层级结构。一张 JPG 并不等于一个
独立患者，也不一定等于一个独立生物样本。图~\ref{fig:medical-hierarchy}给出了本项目最重要的空间层级。

\begin{figure}[H]
\centering
\begin{tikzpicture}[
  node distance=0.34cm,
  box/.style={draw,rounded corners=2pt,minimum height=1.05cm,text width=2.55cm,align=center,fill=CampGray},
  arr/.style={-{Latex[length=2.2mm]},thick,CampBlue}
]
\node[box,fill=CampBlue!11] (organ) {组织/器官或切片\\\small 宏观实验来源};
\node[box,right=of organ,fill=CampGreen!11] (roi) {ROI\\\small 一块连续区域};
\node[box,right=of roi,fill=CampOrange!12] (patch) {patch\\\small $256\times256$ 图块};
\node[box,right=of patch,fill=CampPurple!11] (cell) {细胞/细胞核\\\small 生物学对象};
\node[box,right=of cell,fill=CampRed!10] (pixel) {像素\\\small 一个位置的强度};
\draw[arr] (organ) -- node[above,font=\scriptsize]{选区} (roi);
\draw[arr] (roi) -- node[above,font=\scriptsize]{网格裁切} (patch);
\draw[arr] (patch) -- node[above,font=\scriptsize]{分割/检测} (cell);
\draw[arr] (cell) -- node[above,font=\scriptsize]{采样} (pixel);
\end{tikzpicture}
\caption{医学显微图像的五个层级。左侧更接近独立生物/实验单位，右侧更接近神经网络直接处理的数值。}
\label{fig:medical-hierarchy}
\end{figure}

\begin{itemize}
  \item \textbf{组织/切片层：}实验者从某个生物样本得到组织切片。一张切片可包含大量细胞和多种组织结构。
  当前公开目录没有患者或玻片标识，manifest 的 \code{organ} 也为 \code{unknown}，所以本报告不能把
  “一个 ROI”冒充“一个患者”。
  \item \textbf{ROI 层：}ROI 是切片中一块连续的感兴趣区域。不同 ROI 通常比同一 ROI 内的两个 patch
  更接近统计独立，因此训练/验证划分必须优先按 ROI 分组。
  \item \textbf{patch 层：}为了放入 GPU，大区域被切成 $256\times256$ 小图。文件名里的行、列坐标描述
  它在 ROI 网格中的位置；相邻 patch 可能共享相似的组织纹理和细胞群。
  \item \textbf{细胞层：}病理学家关心的对象往往是细胞，例如“有多少 CD68 相关阳性细胞”“信号是否围绕
  某些细胞核出现”。要从图像得到细胞级量，通常还要先分割细胞核或细胞区域。
  \item \textbf{像素层：}神经网络最终输出每个位置的亮度。像素是采样值，不是一个细胞，也不是一个分子。
  一枚细胞核通常覆盖许多像素，一个像素也可能混合多个亚细胞结构和背景光。
\end{itemize}

这解释了“像素级指标很好但细胞级结论可能错误”：模型可能把一个真实的亮点轻微平移或摊成一片较弱
的光。大部分像素的误差仍然不大，但被判为阳性的细胞数、阳性细胞位置或局部表达强度已经发生变化。

\subsection{一个细胞里有哪些结构，marker 为什么出现在不同位置}

理解 marker 前，先把细胞极度简化为四个区域：\textbf{细胞核}储存 DNA；\textbf{细胞质}包含细胞器和
骨架；\textbf{细胞膜}把细胞与外界分隔并承载许多受体；\textbf{细胞外基质}位于细胞之间并支撑组织。
真实细胞没有规则圆形边界，而且二维切片只截取三维细胞的一层。图~\ref{fig:cell-marker-map}只是空间
位置的教学示意，不代表数据中的真实颜色、比例或细胞类型。

\begin{figure}[H]
\centering
\resizebox{0.96\textwidth}{!}{%
\begin{tikzpicture}[scale=0.82]
  \fill[CampGreen!8] (-6,-2.5) rectangle (6,2.5);
  \draw[CampGreen!55,dashed,thick] (-6,-2.5) rectangle (6,2.5);
  \fill[CampOrange!12] (0,0) ellipse (3.55 and 2.05);
  \draw[CampOrange,very thick] (0,0) ellipse (3.55 and 2.05);
  \fill[CampBlue!32] (-0.75,0.15) ellipse (1.2 and 0.92);
  \draw[CampBlue,thick] (-0.75,0.15) ellipse (1.2 and 0.92);
  \foreach \x/\y in {1.25/0.8,1.7/0.2,1.15/-0.55,2.15/-0.65,0.55/1.15}
    {\fill[CampRed!70] (\x,\y) circle (0.12);}
  \draw[-{Latex},thick,CampBlue] (-4.8,1.75) -- (-1.45,0.75);
  \node[align=right,anchor=east] at (-4.9,1.75) {DAPI\\\small 主要标出 DNA/细胞核};
  \draw[-{Latex},thick,CampRed] (4.95,1.45) -- (1.75,0.55);
  \node[align=left,anchor=west] at (5.0,1.45) {CD68\\\small 常见胞质颗粒样信号};
  \draw[-{Latex},thick,CampPurple] (5.0,-0.2) -- (3.45,-0.35);
  \node[align=left,anchor=west] at (5.05,-0.2) {HLA-DR/CD45RO\\\small 膜相关，解释依赖细胞类型};
  \draw[-{Latex},thick,CampGreen!80!black] (-5.0,-1.65) -- (-2.7,-1.4);
  \node[align=right,anchor=east] at (-5.05,-1.65) {Vimentin\\\small 胞质中间丝/间质相关模式};
  \node at (0,-2.15) {同一位置的不同通道在回答不同的分子问题};
\end{tikzpicture}
}
\caption{五类通道的典型亚细胞位置示意。marker 是分子探针，不是严格一一对应的细胞类别。}
\label{fig:cell-marker-map}
\end{figure}

\subsection{本项目五个通道的医学含义}

训练目录中同一个 basename 在五个 marker 文件夹各有一张图。它们可看作同一空间位置的五种“观察
方式”。表~\ref{tab:marker-primer}中的“常见关联”用于建立直觉，不应被当成临床诊断规则。

\begin{longtable}{>{\raggedright\arraybackslash}p{0.103\textwidth}
                    >{\raggedright\arraybackslash}p{0.165\textwidth}
                    >{\raggedright\arraybackslash}p{0.195\textwidth}
                    >{\raggedright\arraybackslash}p{0.18\textwidth}
                    >{\raggedright\arraybackslash}p{0.162\textwidth}}
\caption{本项目五个通道的入门解释}\label{tab:marker-primer}\\
\toprule
通道 & 主要分子/结构 & 常见医学关联 & 图像中可期待的空间形态 & 不能直接推出什么 \\
\midrule
\endfirsthead
\toprule
通道 & 主要分子/结构 & 常见医学关联 & 图像中可期待的空间形态 & 不能直接推出什么 \\
\midrule
\endhead
DAPI & 与双链 DNA 结合的核染料 & 提供细胞核位置、大小、形状和局部核密度 & 一个个核样亮区域；不同细胞的核形态可不同 & 看不见完整细胞膜；不能仅凭核图确定某蛋白一定表达 \\
CD68 & 内体/溶酶体相关糖蛋白 & 常用于辅助观察单核细胞/巨噬细胞相关信号 & 常见胞质内、核周或颗粒样信号，而不是填满细胞核 & 不是巨噬细胞的绝对专属身份证；阳性不自动等于疾病诊断 \\
CD45RO & CD45/PTPRC 的一种异构体 & 常与人类记忆/活化 T 细胞表型相关，但不同免疫亚群需联合其他 marker 判断\cite{memory_t_cells} & 作为白细胞表面蛋白，常呈细胞轮廓/膜相关分布 & 不能只用该通道确定所有 T 细胞亚型、抗原特异性或功能状态 \\
HLA-DR & 人类 MHC II 类分子 & 抗原呈递细胞表面常见，也可在炎症刺激下于其他细胞上调\cite{hla_dr_review} & 以细胞膜相关信号为主，也受胞内运输影响 & 不是单一细胞类型专属；强信号不能直接等同某种疾病 \\
Vimentin & III 型中间丝细胞骨架蛋白 & 常见于间充质/基质相关细胞，且在多种细胞环境中表达 & 细胞质内纤维状、网状或弥散结构；不与细胞核完全重合\cite{hpa_vimentin} & 不能把所有 Vimentin 阳性区域都判成同一种细胞或固定生物过程 \\
\bottomrule
\end{longtable}

因此，从 DAPI 预测 CD68 等通道并不是简单地给细胞核“换颜色”。例如，一个 DAPI 核周围是否出现 CD68
颗粒，与细胞身份和状态有关；而 DAPI 本身没有直接测量该蛋白。网络只能利用训练集中反复出现的核形态、
密度、组织纹理和蛋白信号之间的\textbf{统计相关性}。

\subsection{把医学概念放回一张真实的五通道 patch}

前面的示意图只说明“marker 可能出现在哪里”。图~\ref{fig:real-five-channel}来自本地官方训练集的
同一个真实空间位置 \code{ROI007_09_11}，能更直观地说明五个通道并不是五张无关照片。前五格直接
显示原始 uint8 灰度，没有对每个通道单独拉伸对比度；右下角只是教学用伪彩叠加，把 DAPI 映射为蓝色、
CD68 映射为红色。数据文件本身并没有这些蓝红颜色。

\begin{figure}[H]
\centering
\includegraphics[width=0.92\textwidth]{docs/essay/figures/marker_channels_roi007_09_11.png}
\caption{官方训练集同一位置的五个真实 marker 通道及 DAPI--CD68 教学伪彩叠加。}
\label{fig:real-five-channel}
\end{figure}

读图时可以按下面的顺序观察：

\begin{enumerate}
  \item \textbf{先看 DAPI：}一个个椭圆或不规则的亮区主要对应细胞核。它能提供核的位置和密度，
  但没有画出完整细胞边界。
  \item \textbf{再横向比较 marker：}HLA-DR 在图中部有较集中的亮结构，CD45RO 的强度和纹理又不同；
  Vimentin 与 CD68 在不少位置呈胞质/轮廓相关信号。它们共享空间坐标，却不是 DAPI 的复制品。
  \item \textbf{最后看叠加图：}蓝色核与红色 CD68 信号经常相邻、包围或部分重叠，而不是严格像素重合。
  这正是“核层级信息”和“胞质/蛋白层级信息”的区别。
\end{enumerate}

这些观察只能描述本 patch 的图像形态，不能据此给每个细胞下临床诊断。例如，不能把所有红色区域直接
数成巨噬细胞，也不能因为某个核周围没有明显红色就断言该细胞绝对不表达 CD68。严谨的细胞级结论还需
细胞分割、多 marker 联合判定、实验阈值和病理专家复核。

图~\ref{fig:real-cell-zoom}进一步放大同一 patch 中 CD68 平均强度较高的 $96\times96$ 窗口。红框和
裁剪位置由脚本根据 CD68 局部均值自动选择，并不是人工挑出的“典型阳性细胞”或医学标注。

\begin{figure}[H]
\centering
\includegraphics[width=0.90\textwidth]{docs/essay/figures/dapi_cd68_cell_zoom.png}
\caption{同一真实 patch 的全图与局部放大。DAPI 主要显示核，CD68 信号更多位于核周围和胞质区域。}
\label{fig:real-cell-zoom}
\end{figure}

局部图也解释了为什么模型难：DAPI 给出蓝色核的形状，却没有直接告诉网络红色信号应当沿哪个核的哪一侧
增强。网络只能从训练样本中学习常见条件关系；当 DAPI 外观相似而 CD68 状态不同时，像素回归会倾向输出
较平滑的平均结果。

\subsection{本项目的 label 到底是什么}

计算机课程里，label 常指“猫/狗”或 $0/1$。本任务不同：主要监督信号是一整张目标 marker 图，称为
\term{稠密标签}{dense target}。每个目标像素都是一个回归值，而不是给整个 patch 一个类别。

\begin{table}[H]
\centering
\caption{数据样本中的字段与“标签”含义}
\label{tab:data-label-semantics}
\begin{tabularx}{\textwidth}{>{\raggedright\arraybackslash}p{0.18\textwidth}
                              >{\raggedright\arraybackslash}p{0.22\textwidth}X}
\toprule
信息 & 计算机表示 & 含义与用途 \\
\midrule
DAPI 输入 & $H\times W\times C$ 图像 & 模型可见的条件信息；当前 JPG 存储为 RGB，但三通道相同。\\
目标图 & 同尺寸连续强度图 & train 中 HLA-DR、CD45RO、Vimentin、CD68 都存在；当前初赛配置只把 CD68 作为训练目标。\\
marker 名 & 字符串/任务 ID & 告诉多任务模型要预测哪种蛋白通道；它不是患者诊断。\\
ROI、row、col & 分组与整数坐标 & 用于防止泄漏、恢复邻域和按区域汇总；不直接参与当前默认 context-disabled 模型。\\
split & train/val/test & train 用于更新权重，val 用于选模型，test 只有 DAPI 且隐藏目标由竞赛平台保管。\\
organ & 元数据字符串 & 当前本地 manifest 为 \code{unknown}；在没有官方元数据时不能从文件夹结构臆造器官标签。\\
\bottomrule
\end{tabularx}
\end{table}

本地 official 数据的结构是：训练集中每个 marker 各 6296 张、可组成 6296 组同名配对；按真实 ROI
分组后得到 5004 个 train patch 和 1292 个 val patch。测试目录有 1346 张 DAPI，但没有任何目标图。
所以 \textbf{val 分数是本地可计算的开发指标，test 分数必须把预测提交到平台后才能得到}。

还要区分“图像强度”和“蛋白分子数”。荧光图中的亮度受染色效率、曝光、探测器、背景、自发荧光、
漂白和饱和等因素影响；若没有逐批校准，像素值通常只是相对强度，不等于一个可以跨实验直接比较的
绝对分子计数\cite{fluorescence_quantitation}。因此网络拟合的是数据制作流程下的目标图像分布，而不是
把 $0\sim255$ 机械地翻译成“多少个蛋白分子”。

\subsection{配对与配准：同名文件为什么重要}

监督学习需要知道哪一张 DAPI 对应哪一张目标图。本工程用 marker 文件夹和规范化后的同名 stem 建立
配对，例如：

\begin{lstlisting}[style=shell]
train/DAPI/ROI000_03_12.jpg
train/CD68/ROI000_03_12.jpg
\end{lstlisting}

相同的 \code{ROI000_03_12}表示二者应来自同一网格位置；这叫\term{空间配对}{spatial pairing}。
但“文件名相同”不等于所有边缘都精确重合：制样、通道采集、光学系统和压缩仍可能造成轻微位移或噪声。
本地审计得到目标通道与 DAPI 的中位平移为 0 像素，说明没有发现系统性整体错位，所以默认不做配准；
这不代表每个细胞的所有亚细胞信号都应该与 DAPI 核重合。CD68 本来就主要不是核内信号。

若训练/验证把同一 ROI 的相邻 patch 随机拆开，模型可能在验证时看见几乎相同的组织环境，导致分数虚高。
这叫\term{数据泄漏}{data leakage}。本项目按 ROI 分组，当前 train 与 val 的 ROI 不重叠。对于未来含患者或
玻片 ID 的数据，应该优先按患者/玻片分组，因为它们比 ROI 更高一层。

\subsection{从光子到 Tensor：计算机实际接收了什么}

荧光显微镜把某个波段的光转换成探测器读数，保存前又可能经过曝光、增益、位深转换和压缩。进入当前
工程后，一张图经历下面的表示变化：

\begin{enumerate}
  \item 磁盘上的 JPEG 解码为 $256\times256\times3$ 的 uint8 数组，元素为整数 $0\sim255$；
  \item 图像库按 RGB 顺序读取。审计发现三通道逐像素相同，所以医学信息上是一个灰度强度通道；
  \item 数据集把像素除以 255，变成 float32 的 $0\sim1$；
  \item 维度从 HWC 调整为 PyTorch 的 CHW，再加 batch 维成为 $B\times C\times H\times W$；
  \item 网络输出同范围的浮点预测；提交时先四舍五入量化回 uint8，再以固定 JPEG 参数编码。
\end{enumerate}

JPEG 是有损格式。同一个浮点预测在“内存 float”“量化 uint8”“写成 JPG 后重新读取”三个域里并不完全
一致，这也是项目同时报告三域指标的原因。背景、自发荧光和探测器噪声会减少有效动态范围，饱和像素则
丢失超过上限的强度信息；这些都是定量荧光显微分析的常见限制\cite{fluorescence_quantitation}。

\subsection{模型学的是条件预测，不是把不可见信息恢复出来}

更准确的数学说法是，训练数据从某个联合分布产生 $(X_{\mathrm{DAPI}},Y_{\mathrm{marker}})$，模型试图学习

\begin{equation}
  \widehat{Y}=f_\theta(X)\approx \mathbb{E}[Y\mid X].
\end{equation}

如果两个细胞的 DAPI 外观几乎相同，但一个表达 CD68、另一个不表达，那么仅凭 DAPI 无法唯一判别。
在均方误差等损失下，模型容易给出两种可能性的平均，即一个偏弱、偏平滑的信号。这不是简单增加网络
参数就一定能解决的问题，而是\term{条件不确定性}{conditional uncertainty}：输入中可能缺少决定目标的
分子信息。

因此应避免三种过度表述：其一，“模型检测到了 CD68 蛋白”；准确说法是“模型预测了 CD68 图像”；
其二，“亮点就是巨噬细胞”；准确说法是“亮点与 CD68 相关，细胞身份仍需联合形态和其他 marker”；
其三，“SSIM 提高就有临床价值”；临床价值还需要外部队列、细胞级分析和专家验证。

\subsection{像素级、细胞级、patch 级与 ROI 级评价}

表~\ref{tab:evaluation-levels}把后文经常出现的“层级”拆开。一个指标属于哪一层，决定了它能回答什么问题。

\begin{longtable}{>{\raggedright\arraybackslash}p{0.105\textwidth}
                    >{\raggedright\arraybackslash}p{0.22\textwidth}
                    >{\raggedright\arraybackslash}p{0.25\textwidth}
                    >{\raggedright\arraybackslash}p{0.27\textwidth}}
\caption{不同评价层级的含义}\label{tab:evaluation-levels}\\
\toprule
层级 & 计算对象 & 可以回答的问题 & 局限与常见指标 \\
\midrule
\endfirsthead
\toprule
层级 & 计算对象 & 可以回答的问题 & 局限与常见指标 \\
\midrule
\endhead
像素级 & 每个坐标的预测强度与真值强度 & 整体灰度和局部结构有多接近？ & MSE、PSNR、SSIM；不理解细胞身份，且大量暗背景会稀释稀有亮点的错误。\\
细胞级 & 先用 DAPI 分割核/细胞，再对每个对象汇总 marker & 阳性细胞比例、单细胞平均/最大强度、信号与核的距离是否合理？ & 依赖分割和阳性阈值；本项目当前标准 Validator 尚未计算。\\
patch 级 & 一张 $256\times256$ 图的指标或统计量 & 哪些局部图块最难？性能是否随亮度、密度变化？ & 同一 ROI 内 patch 高度相关，不能把它们全当独立样本。\\
ROI 级 & 同一 ROI 的全部 patch & 一个连续组织区域的平均表现和空间一致性如何？ & ROI 数通常远小于 patch 数；本项目用 ROI 聚合和成对 bootstrap 做可靠比较。\\
切片/患者级 & 同一玻片或患者的多个 ROI & 方法能否跨患者泛化，是否改变患者级结论？ & 当前数据没有可验证的患者/玻片 ID，不能声称通过患者级验证。\\
\bottomrule
\end{longtable}

一个简单反例能说明为什么要分层：真值里有 2 个强阳性细胞，预测把相同的总亮度均匀分给 10 个细胞。
两张图的总强度可能相近，某些像素平均误差也未必极端，但“阳性细胞数量”和“信号集中在哪些细胞”已经
完全不同。反过来，如果预测只比真值平移 1 个像素，细胞级阳性结论可能不变，像素级 MSE 却会明显增大。

\subsection{SSIM、PSNR 与医学意义应该怎样对应}

给真值图 $Y$ 和预测图 $\widehat{Y}$，均方误差为

\begin{equation}
  \operatorname{MSE}=\frac{1}{N}\sum_{i=1}^{N}(Y_i-\widehat{Y}_i)^2,
\end{equation}

当图像范围归一化为 $[0,1]$ 时，PSNR 为

\begin{equation}
  \operatorname{PSNR}=10\log_{10}\!\left(\frac{1}{\operatorname{MSE}}\right).
\end{equation}

PSNR 高表示平均平方误差小，但它对整体亮度范围敏感，也不知道一个误差落在背景还是关键阳性细胞上。
SSIM 在局部窗口比较亮度、对比度和结构\cite{ssim}，比单纯 MSE 更接近“纹理是否相似”，但它也没有
细胞语义：两团形状相似却属于不同细胞的信号，SSIM 不会自动理解其医学差别。

本报告后文的 float/uint8/JPG 指标仍属于\textbf{同一套像素/局部结构指标在三种文件表示上的计算}，
不是三个医学终点。若要进一步研究医学一致性，可在不改变竞赛主指标的前提下增加：单细胞 marker
强度相关、阳性细胞比例误差、细胞间强度排序的 Spearman 相关、ROI 内空间分布差异等。它们必须固定
细胞分割方法和阈值，并报告失败样本，不能挑一张好看的热力图代替统计检验。

\subsection{读后续网络章节前应掌握的最小知识清单}

\begin{enumerate}
  \item 一张 target 图是\textbf{稠密回归标签}，不是一个类别编号；
  \item DAPI 主要描述细胞核，CD68/HLA-DR/CD45RO/Vimentin 描述的是不同分子和亚细胞位置；
  \item marker 与细胞类型是“相关但不绝对唯一”的关系，细胞身份通常需要多 marker 与形态联合判断；
  \item 一个 ROI 含很多相邻 patch，随机按 patch 划分会造成泄漏；
  \item 像素、细胞、patch、ROI、患者是不同统计单位，不能混用样本量；
  \item 荧光亮度是受实验流程影响的相对测量，不应直接称为蛋白分子数；
  \item 网络预测条件分布中的常见结果，不能凭 DAPI 创造输入中完全不存在的确定信息；
  \item PSNR/SSIM 衡量图像相似度，不自动证明细胞级或临床级正确；
  \item test 没有公开目标，验证命令只能评估 val，平台分数才是隐藏 test 的结果；
  \item JPEG 量化/压缩会改变像素，所以最终竞赛比较应重点看 JPG round-trip 指标和平台结果。
\end{enumerate}

\section{命令行入门：为什么用 CLI 而不是 \texttt{python train.py}}

\subsection{CLI 是什么}

\term{CLI}{Command Line Interface，命令行接口}是通过输入文字命令控制程序的方式。图形界面像餐厅
菜单，CLI 更像把“菜名、份量、备注”写成一行交给厨房。它看起来不如按钮直观，但非常适合服务器：
命令可以复制、记录、重跑，也能清楚说明使用了哪个配置和数据目录。

本项目推荐的训练命令是：

\begin{lstlisting}[style=shell]
python -m virtual_staining.cli --log-root log autodl-run \
  --config configs/initial_round_cd68_max_v4.yaml \
  --data-root AUTO --target CD68 --run-id cd68_max_v4_seed2026
\end{lstlisting}

\subsection{把命令逐段拆开}

\begin{longtable}{>{\raggedright\arraybackslash}p{0.27\textwidth}p{0.66\textwidth}}
\toprule
命令片段 & 含义 \\
\midrule
\endfirsthead
\toprule
命令片段 & 含义 \\
\midrule
\endhead
\code{python} & 启动当前环境中的 Python 解释器。AutoDL 上应先激活装好 PyTorch 的环境。\\
\code{-m} & 告诉 Python：“后面给的是可导入模块名，请按模块运行”，而不是磁盘上的单个脚本路径。\\
\code{virtual_staining.cli} & Python 模块路径，对应源码 \code{src/virtual_staining/cli.py}。安装项目后，Python 能从 \code{src/} 布局中找到它。\\
\code{--log-root log} & 全局选项。每次 CLI 调用会在根目录 \code{log/} 下创建轻量日志目录，记录环境、命令、配置、epoch 指标和错误；成功结束后还会打成日志 ZIP。该参数必须放在子命令之前。\\
\code{autodl-run} & 子命令，表示执行“环境检查、数据准备、训练和验证”的一键编排。\\
\code{--config ...yaml} & 指定本次使用的 YAML 配置。YAML 就像一张集中填写超参数的表。\\
\code{--data-root AUTO} & 不硬编码目录，让发现程序在受限候选路径中评分并选择数据根；当前应发现 \code{dataset/official}。\\
\code{--target CD68} & 通用训练/验证命令可用的目标参数。需要注意：当前 \code{autodl-run} 包装器实际从 YAML 的 \code{data.targets} 和 \code{data.submit_targets} 读取唯一目标；本配置本来就是 CD68，因此该参数起到说明意图的作用，但不要误以为它会覆盖包装器中的 YAML 目标。\\
\code{--run-id ...} & 本次实验的唯一名字，同时用于输出目录和日志目录。已有非空输出目录不会被静默覆盖。\\
\bottomrule
\end{longtable}

\subsection{为什么不是直接运行 \texttt{train.py}}

项目根目录并不存在一个承担全部工作的 \code{train.py}。代码采用标准
\term{src 布局}{src layout}：可导入包位于 \code{src/virtual_staining/}。使用 \code{-m} 有三点好处：

\begin{enumerate}
  \item Python 按包规则解析模块间导入，不依赖你碰巧站在哪个子目录；
  \item 所有功能共享同一个入口，却可通过子命令区分 \code{train}、\code{validate}、\code{predict}、\code{autodl-submit} 等流程；
  \item 参数解析、日志、异常退出码和配置合并方式保持一致，便于服务器自动化。
\end{enumerate}

如果已经执行 \code{pip install -e .}，包会以 editable（可编辑）方式注册到当前 Python 环境；修改
\code{src/} 后不需要反复复制包。若不安装，也可以把根目录下的 \code{src} 加进 \code{PYTHONPATH}。

\subsection{从入口到训练器：\texttt{autodl-run} 内部流程}

图~\ref{fig:cli-flow}展示了命令的路由关系。

\begin{figure}[H]
\centering
\resizebox{0.98\textwidth}{!}{%
\begin{tikzpicture}[
  node distance=8mm and 9mm,
  box/.style={draw,rounded corners,fill=CampGray,minimum height=9mm,align=center,text width=31mm},
  action/.style={draw,rounded corners,fill=CampBlue!10,minimum height=9mm,align=center,text width=34mm},
  arr/.style={-{Latex[length=2.2mm]},thick}
]
\node[box] (python) {\code{python -m}\\加载模块};
\node[box,right=of python] (main) {\code{main()}\\程序入口};
\node[box,right=of main] (parser) {\code{build_parser()}\\注册并解析子命令};
\node[action,right=of parser] (route) {autodl-run\\一键编排};
\draw[arr] (python)--(main);
\draw[arr] (main)--(parser);
\draw[arr] (parser)--(route);
\node[action,below left=13mm and 9mm of route] (data) {环境报告、数据发现、\\manifest、数据审计};
\node[action,below=13mm of route] (train) {\code{_run_train()}\\模型与 Trainer.fit};
\node[action,below right=13mm and 9mm of route] (val) {validate handler\\验证并写报告};
\draw[arr] (route)--(data);
\draw[arr] (data)--(train);
\draw[arr] (train)--(val);
\end{tikzpicture}}
\caption{\texttt{autodl-run} 的命令解析与主要调用链}
\label{fig:cli-flow}
\end{figure}

真实代码的执行顺序如下：

\begin{enumerate}
  \item 模块末尾调用 \code{main()}；\code{build_parser()} 创建 \code{argparse} 解析器并注册所有子命令。
  \item 参数解析完成后，\code{ExperimentLogSession.from_args()} 先在 \code{log/<run-id>/} 创建 invocation、environment 和 status 文件，并把控制台输出同时写入日志。
  \item 解析器发现子命令是 \code{autodl-run}，调用 \code{command_autodl_run()}。
  \item \code{command_env()} 收集 Python、PyTorch、CUDA、GPU、CPU 和内存信息，写入 \code{artifacts/environment.json}。
  \item \code{discover_data_root()} 对有限候选目录打分；当前数据根为 \code{dataset/official}。
  \item \code{_ensure_manifests()} 配对文件并建立 train、val、test 与 smoke CSV；随后 \code{audit_data()} 统计尺寸、通道、值域、相关性和对齐情况。
  \item \code{_run_train()} 合并配置、构建带 hard-mix activity sampler 的 DataLoader、\code{MultiMarkerRestorer}、scheduled composite loss、分组 fused AdamW、Trainer 与 Validator；\code{Trainer.fit()} 最多训练 100 epoch，并受 early stopping 与 8.5 小时预测式安全停点约束。
  \item v4 每 5 epoch 比较 raw 与 EMA；15\% batch 还会执行一次随机 D4 等变一致性分支。训练结束后，\code{command_validate()} 加载最佳 checkpoint 中记录的选中权重来源，输出逐图和汇总指标；当前 v4 推理固定使用 \code{tta=d4} 与 \code{tta\_policy=configured}。
  \item \code{pipeline_report.json} 汇总环境、数据、训练和验证路径。CLI 成功结束时，日志会复制关键轻量文件并生成 \code{autodl-run_log_bundle.zip}。
\end{enumerate}

\subsection{\texttt{autodl-submit} 内部流程}

提交命令不会再训练。它依次做以下事情：

\begin{enumerate}
  \item 若给了 \code{--checkpoint} 就使用它；否则在 \code{outputs/**/best_ssim.ckpt} 中选修改时间最新的文件。
  \item 重新发现数据并建立 test manifest，确认 test 不是空表。
  \item \code{command_predict()} 根据 checkpoint 内保存的模型架构和 ImageSpec 构建模型，选择 raw 或 EMA 权重，按需要执行 D4 TTA，并把预测保存到 \code{predictions/CD68/}。
  \item \code{command_make_submission()} 将预测复制为规定命名，构建 \code{results/test/CD68/} 并压缩为 ZIP。已经是 JPEG 的预测会直接复制，避免二次有损编码。
  \item \code{validate_submission()} 检查文件名、数量、尺寸、模式、像素范围、纯黑/纯白异常、ZIP 层级和 CRC。任一硬错误都会令命令返回非零退出码。
\end{enumerate}

\section{数据集格式与加载}

\subsection{本地官方数据的真实目录与数量}

当前工作区的真实目录为：

\begin{lstlisting}[style=shell]
dataset/official/
|-- train/
|   |-- DAPI/       6296 JPG
|   |-- HLA-DR/     6296 JPG
|   |-- CD45RO/     6296 JPG
|   |-- Vimentin/   6296 JPG
|   `-- CD68/       6296 JPG
`-- test/
    `-- DAPI/       1346 JPG
\end{lstlisting}

训练目录每个 marker 的文件名一一对应，总计 6296 组配对样本。当前初赛配置只选择 CD68，因此每次
训练实际读取 DAPI 与 CD68；其余三个目标仍在磁盘和 manifest 中，但不参与本次单目标 loss。
test 只有 1346 张 DAPI，没有 CD68 标签，所以本地无法为 test 计算 SSIM 或 PSNR。

\subsection{真实样本显示了哪些数据分布特征}

仅展示一张“好看的图”容易造成误解。图~\ref{fig:cd68-activity-quantiles}先对 6296 张训练 CD68 图计算
归一化平均强度，再选取接近第 10、50、90 百分位的三个样本。每行左侧是模型输入 DAPI，中间是训练
目标 CD68，右侧是教学伪彩叠加；\code{mean}是 CD68 全图平均强度，\code{>0.25}表示强度超过 0.25
的像素比例。

\begin{figure}[H]
\centering
\includegraphics[width=0.92\textwidth]{docs/essay/figures/cd68_activity_quantiles.png}
\caption{按 CD68 全图平均强度选取的低、中、高活动度真实配对样本。这里的“活动度”是计算机统计量，不是临床阳性分级。}
\label{fig:cd68-activity-quantiles}
\end{figure}

这三行揭示了四个与建模直接相关的特征：

\begin{itemize}
  \item \textbf{目标强度分布不均匀：}不同 patch 的背景、亮区面积和局部峰值相差明显。若训练只优化
  全图平均误差，大量暗背景可能压过少量但重要的亮结构。
  \item \textbf{DAPI 与 CD68 不是确定映射：}核密度、核亮度与 CD68 有统计联系，但从三行不能得到
  “DAPI 越亮，CD68 必然越亮”的单调规则。
  \item \textbf{标签是连续图而非二分类 mask：}CD68 目标包含从暗背景到强信号的连续值。报告中的
  0.25 只用于展示像素占比，不是项目训练标签，也没有被当作临床阳性阈值。
  \item \textbf{结构与强度都要学习：}网络既要预测哪里出现信号，又要估计信号有多强；这也是复合损失
  同时使用 MSE/Charbonnier 与 SSIM、梯度、金字塔项的动机。
\end{itemize}

选择过程可由 \code{docs/essay/build_report_figures.py} 完整复现，具体样本、统计量和显示规则记录在
\code{docs/essay/figures/report_figure_provenance.json}。脚本只读取原始数据，不会覆盖、重编码或移动
\code{dataset/official} 中的任何文件。

\subsection{“RGB 文件”和“灰度内容”不是同一概念}

本地 \code{artifacts/data_audit.json} 的全量审计结果为：所有 DAPI 和四个目标图都是
$256\times256$、8 位 JPEG、PIL mode \code{RGB}、3 个存储通道；同时，三个通道之间的平均差和
最大差均为 0。因此，图像\textbf{内容}在逻辑上是灰度，\textbf{文件容器}却是 RGB。

\begin{table}[H]
\centering
\caption{本地官方数据的存储与逻辑通道事实}
\begin{tabular}{lrrrrl}
\toprule
Marker & 图数 & 尺寸 & 存储通道 & 逻辑通道 & PIL mode \\
\midrule
DAPI & 6296 & $256\times256$ & 3 & 1 & RGB \\
HLA-DR & 6296 & $256\times256$ & 3 & 1 & RGB \\
CD45RO & 6296 & $256\times256$ & 3 & 1 & RGB \\
Vimentin & 6296 & $256\times256$ & 3 & 1 & RGB \\
CD68 & 6296 & $256\times256$ & 3 & 1 & RGB \\
\bottomrule
\end{tabular}
\label{tab:image-spec}
\end{table}

\important{\code{initial_round_cd68_max_v4.yaml} 明确写入 1 个逻辑输入/输出通道，配置加载器还会合并
\code{configs/resolved_local.yaml}；该文件根据审计结果设置 \code{data.input_channels: 1}，并把四个
target 的逻辑通道数都设为 1。因此当前 Dataset 与模型实际按 $[B,1,256,256]$ 运行，输出也是单通道
逻辑张量。这里的“单通道训练”不等于把原始文件改成 mode L：官方 JPEG 仍保持只读 RGB，提交保存时
仍按 checkpoint 中记录的 RGB storage mode 恢复三个相同通道。}

\subsection{文件名、ROI 与二维坐标}

文件名使用 \code{ROI{编号}_{行}_{列}.jpg}。manifest 中旧的 \code{parse_roi_id()} 会提取
\code{ROI000} 用于分组；更严格的 \code{parse_patch_coordinate()} 可进一步解析行列，供网格审计、
边界/内部分类和邻域模型使用。当前 v4 配置的 \code{model.context.enabled=true}，但
\code{require_verified_grid=true}：只有文件名坐标、行列方向、边界连续性和 split 隔离审计全部通过，
才会读取同一 ROI 的 $3\times3$ DAPI 邻域；缺失位置由 mask 表示，目标 marker 永远不作为 context 输入。

坐标有空洞并不一定代表数据损坏。例如 ROI 网格可能只保留组织区域，背景 patch 被排除。程序只把
真实存在的文件写入 manifest，不会凭空生成缺失坐标。

\subsection{manifest 是什么}

\term{manifest}{清单文件}是 CSV 表格。它不存像素，只存“到哪里找图”和“这张图属于谁”。每一行
包含 organ、split、ROI、patch ID、规范化 key、DAPI/目标相对路径、宽高、通道、mode、文件大小、
SHA-1 指纹、是否配对及缺失标记等。这样，训练过程可以被审计：某个预测出错时，能追溯到原始文件。

路径写成相对于数据根的形式，例如
\code{train/DAPI/ROI000_00_00.jpg}，运行时再与数据根拼接。代码使用 \code{pathlib.Path}，因此支持
中文和空格路径。

\subsection{按 ROI 划分训练集与验证集}

官方目录没有单独的 val 文件夹。程序先从文件名得到 25 个训练 ROI（ROI000--ROI024），再以种子
2026 按组抽取约 20\% ROI 作为本地 val。当前清单结果见表~\ref{tab:split}。

\begin{table}[H]
\centering
\caption{当前 manifest 划分结果}
\begin{tabular}{lrrp{0.47\textwidth}}
\toprule
集合 & patch 数 & ROI 数 & 说明 \\
\midrule
train & 5004 & 20 & 用于反向传播与参数更新 \\
val & 1292 & 5 & ROI004、ROI005、ROI006、ROI011、ROI023；不参与参数更新 \\
test & 1346 & 5 & 只有 DAPI；用于生成提交 \\
\bottomrule
\end{tabular}
\label{tab:split}
\end{table}

当前泄漏报告显示 train 与 val 的 canonical key、DAPI 哈希和 group ID 交集均为空，ROI 来源为可靠的
文件名正则，因而这次 ROI 分组划分是可验证的。为什么不能简单随机抽 20\% patch？因为同一 ROI 的
相邻 patch 很相似。如果左边 patch 进训练集、右边 patch 进验证集，模型可能靠记忆局部区域获得虚高
分数。按 ROI 分组相当于把整本练习册分开，而不是把同一道题的上下半部分分别放到训练和考试中。

\subsection{Dataset 与 DataLoader 如何把 JPG 送入模型}

\code{VirtualStainingDataset.__getitem__()} 的核心步骤是：

\begin{enumerate}
  \item 从 manifest 取一行，拼出 DAPI 和 CD68 的绝对路径；
  \item PIL 解码 JPEG。灰度存储会转 L，其他存储会转 RGB；当前文件因此得到 HWC RGB 数组；
  \item 因逻辑通道数为 1，对三个相同 RGB 通道取均值并四舍五入，再转为 CHW 张量、除以 255，得到 \code{float32} 的 $[0,1]$ 值；
  \item 检查 DAPI 与目标空间尺寸完全相同；
  \item 训练时调用成对增强，最后再次转 \code{float32} 并截断到 $[0,1]$；
  \item 返回输入、targets 字典、stem、ROI、organ、split 和路径元数据。
\end{enumerate}

v4 YAML 在 AutoDL/Linux 上使用 \path{batch_size=6}、\path{gradient_accumulation=2}、\path{num_workers=8}、\path{prefetch_factor=4}、
\path{pin_memory=true} 与 \path{persistent_workers=true}，以并发解码减少 4090 等待。加载器只在 CUDA
可用时真正 pin memory；Windows 本地若 worker 不稳定，可用 CLI 覆盖为 0，此时程序也会关闭 persistent。
每次运行都应以日志中的 effective config 为准，而不能只看源 YAML 猜测最终生效值。

\subsection{PairedTransform 为什么必须“成对”}

若只翻转 DAPI、不翻转 CD68，标签就错位了。因此水平翻转、垂直翻转、$0/90/180/270$ 度旋转和可选
平移会对输入与所有目标同步执行。亮度、对比度、gamma、噪声和模糊则只施加在 DAPI 输入上，用来模拟
采集强度变化，同时保留真实 CD68 标签。当前通过 \code{_make_loader()} 创建普通 PairedTransform 时，
只从配置传入 \code{max_translation}；其余参数使用类的默认值：水平/垂直翻转概率 0.5、旋转概率 1.0、
gamma/亮度/对比度概率各 0.5、噪声概率 0.25、模糊概率 0，平移默认为 0。

\section{模型架构详解}

\subsection{先建立整体直觉：这是一个“先压缩理解，再放大还原”的网络}

当前 v4 候选以 \code{MultiMarkerRestorer} 为核心，可以看成带三条互补特征路线的 U-Net 风格编码器--解码器\cite{unet}。编码器逐级减小宽高、增加通道，
把像素转成更抽象的特征；解码器再逐级放大，并从编码器拿回高分辨率 skip feature，恢复细节。
与普通 U-Net 不同的是，本模型用 NAF block 负责恢复，在 $1/8$ 瓶颈用共享/任务余弦原型混合语义，
并在每个解码尺度用 CD68 task adapter 做条件调制。图~\ref{fig:model-overview}先画出 v4 从稳定 v2
继承的 NAF 主干；图~\ref{fig:v4-triroute}再展开当前真正启用的 Heavy Omni U、自适应频率、内容感知
融合与 ROI context。v3 轻量分支只作为历史对照单独保留，绝不是当前 v4 前向的一部分。

\begin{figure}[H]
\centering
\resizebox{\textwidth}{!}{%
\begin{tikzpicture}[
  node distance=7mm,
  stage/.style={draw,rounded corners,align=center,minimum height=14mm,text width=25mm,fill=CampBlue!10},
  mix/.style={draw,rounded corners,align=center,minimum height=14mm,text width=29mm,fill=CampPurple!12},
  head/.style={draw,rounded corners,align=center,minimum height=14mm,text width=28mm,fill=CampOrange!16},
  io/.style={draw,rounded corners,align=center,minimum height=14mm,text width=22mm,fill=CampGreen!12},
  arr/.style={-{Latex[length=2.2mm]},thick},
  skip/.style={-{Latex[length=2mm]},dashed,thick,CampGreen}
]
\node[io] (in) {DAPI 输入\\$B\times1\times256\times256$\\+ Sobel};
\node[stage,right=of in] (e1) {NAF Stage\\$256^2$，64 通道\\2 blocks};
\node[stage,right=of e1] (e2) {NAF Stage\\$128^2$，128 通道\\2 blocks};
\node[stage,right=of e2] (e4) {NAF Stage\\$64^2$，256 通道\\4 blocks};
\node[stage,right=of e4] (e8) {NAF Stage\\$32^2$，512 通道\\6 blocks};
\draw[arr] (in)--(e1); \draw[arr] (e1)--(e2); \draw[arr] (e2)--(e4); \draw[arr] (e4)--(e8);

\node[mix,below=13mm of e8] (proto) {8 shared + 8 CD68 原型\\cosine attention\\temperature 0.1};
\draw[arr] (e8)--(proto);

\node[stage,below=19mm of e4] (d4) {Decoder + adapter\\$64^2$，256 通道\\2 blocks};
\node[stage,left=of d4] (d2) {Decoder + adapter\\$128^2$，128 通道\\2 blocks};
\node[stage,left=of d2] (d1) {Decoder + adapter\\$256^2$，64 通道\\2 blocks};
\node[head,left=of d1] (outhead) {$1\times1$ head\\sigmoid\\64/128/256 深监督};
\node[io,left=of outhead] (out) {CD68 预测\\$B\times1\times256\times256$};

\draw[arr] (proto)--(d4); \draw[arr] (d4)--(d2); \draw[arr] (d2)--(d1); \draw[arr] (d1)--(outhead); \draw[arr] (outhead)--(out);
\draw[skip] (e4)--(d4);
\draw[skip] (e2) to[bend left=15] (d2);
\draw[skip] (e1) to[bend right=14] (d1);
\end{tikzpicture}}
\caption{v4 继承的 MultiMarkerRestorer NAF 核心主干。图中按实际单逻辑通道标注输入/输出；当前重型三路增量见图~\ref{fig:v4-triroute}。}
\label{fig:model-overview}
\end{figure}

\subsection{历史 v3 对照：轻量 U、固定频率与稳定输出（当前未启用）}

本小节用于解释旧配置和旧 checkpoint，不能拿来描述当前 v4。v3 候选不是把主干替换成一个更大的 Transformer，而是在已验证的 NAF U 形网络旁边增加四条互补、
可由 feature flag 独立关闭的路径，如图~\ref{fig:v3-branches}。这遵循高效卷积复原网络的经验：
局部大核/深度可分离卷积可以用较少参数补充细节\cite{convir,oknet}；频率分解可以显式区分平滑强度与
高频纹理\cite{adair}；训练期等变约束可以减少旋转/翻转后的输出不一致\cite{eqreg}。这些文献只提供
设计依据，本工程是针对 DAPI$\rightarrow$CD68 和单卡 4090 预算的 clean-room 实现，不复制论文权重或结果。

\begin{figure}[H]
\centering
\resizebox{0.98\textwidth}{!}{%
\begin{tikzpicture}[
  node distance=6mm,
  io/.style={draw,rounded corners,align=center,fill=CampGreen!12,text width=24mm,minimum height=12mm},
  core/.style={draw,rounded corners,align=center,fill=CampBlue!10,text width=30mm,minimum height=12mm},
  branch/.style={draw,rounded corners,align=center,fill=CampPurple!12,text width=34mm,minimum height=12mm},
  head/.style={draw,rounded corners,align=center,fill=CampOrange!15,text width=35mm,minimum height=12mm},
  arr/.style={-{Latex[length=2mm]},thick},
  res/.style={-{Latex[length=2mm]},dashed,thick,CampPurple}
]
\node[io] (x) {中心 DAPI\\+ Sobel};
\node[core,right=of x] (enc) {四尺度 NAF\\主编码器};
\node[core,right=of enc] (bottle) {$1/8$ bottleneck\\prototype};
\node[core,right=of bottle] (dec) {共享 decoder\\task adapters};
\node[head,right=of dec] (head) {direct logits\\+ zero-init base/detail};
\node[io,right=of head] (y) {CD68\\预测};
\draw[arr] (x)--(enc)--(bottle)--(dec)--(head)--(y);

\node[branch,above=12mm of enc] (detail) {轻量 Detail U-Net\\大核 depthwise/gated\\注入 $1,1/2,1/4$ skip};
\draw[arr] (x) |- (detail);
\draw[res] (detail)--(enc);

\node[branch,below=13mm of enc] (ctx) {$3\times3$ ROI DAPI tiles\\tiny shared encoder\\FiLM at $1/4,1/8$};
\draw[res] (ctx)--(enc);
\draw[res] (ctx) to[bend right=12] node[below,align=center]{9-token\\cross-attention} (bottle);

\node[branch,below=13mm of bottle] (fft) {parallel spatial/frequency\\NAF + float32 FFT low/high\\channel softmax gate};
\draw[res] (bottle)--(fft);
\draw[res] (fft)--(dec);

\node[branch,below=13mm of head] (d4) {训练期 D4 consistency\\25\% batch 额外一次变换\\推理仍按配置执行 D4};
\draw[res] (d4)--(head);
\end{tikzpicture}}
\caption{历史 v3 增量路径，仅供兼容与对照；当前 v4 已用 Heavy Omni U 和三路门控替换轻量 U/固定频率组合。}
\label{fig:v3-branches}
\end{figure}

\paragraph{轻量 Detail U-Net。}
它使用宽度 $[16,24,32,48]$ 的四尺度 U 形分支，每层由 $7\times7$ depthwise convolution、
$1\times1$ expansion 和乘法 gate 组成。解码结果只向主干的全分辨率、$1/2$、$1/4$ skip 注入；
$1/8$ 留给空间--频率 mixer，避免重复计算。三个注入投影严格零初始化，所以启用分支时第 0 步
不改变旧主干输出。该分支约增加 72,016 个参数和 0.589 GMAC。

\paragraph{空间--频率瓶颈。}
同一个 $32\times32$ bottleneck 一路进入局部 NAF，一路在 float32 中执行 rFFT，使用平滑径向 mask
拆成 low/high bands，再由逐通道 softmax gate 融合。最终 projection 乘零初始化 residual scale 后
接回主干。FFT 强制 float32 是为了 CUDA 数值/shape 兼容，其余卷积继续在 BF16 autocast 下运行。

\paragraph{真实 ROI context。}
文件名 \code{ROI007_09_11} 可解析为 ROI、row、col；数据审计已经确认官方 train/val 没有共享 ROI 或
跨 split 邻接边，行列方向连续。\code{NeighborhoodDataset} 因此可以安全读取中心周围最多 9 张 DAPI，
并返回 valid mask 与二维 offsets。增强会同步变换 tiles、mask 和 offsets；中心 tile 强制与主输入逐位一致。
若换数据后文件名网格无法验证，\code{require_verified_grid=true} 会拒绝训练，不会根据图像内容猜坐标。

\paragraph{稳定 base/detail。}
旧 base/detail 会用随机新头直接替换最终 output head，单次重训风险较高。v3 引入且 v4 继续使用的
\code{base_detail_residual=true} 改为
\begin{equation}
  z_{\mathrm{final}}=z_{\mathrm{direct}}
  +\Delta z_{\mathrm{coarse}}
  +a\tanh(\Delta z_{\mathrm{detail}}),\qquad a=0.5,
\end{equation}
其中 coarse 和 detail 的最终投影都从 0 初始化。因此第 0 步 prediction 与旧 direct head 逐位一致，
但两个投影在第一次 backward 都能获得梯度。

\subsection{v4 重型三路主线：为什么替换轻量 U，却不删除旧模型}

v3 的轻量 Detail U 只有 72,016 个参数，只向三个 skip 输出零初始化更新，
没有独立的解码细化和 $1/8$ 瓶颈路线。它适合做低风险增量，但不能满足“重型 U-Net
与另外两路平行后再融合”的设计意图。因此 v4 使用新的
\path{models/heavy_detail_unet.py}，宽度为 $[56,88,128,176]$，编码深度为
$[3,4,6,6]$，解码细化深度为 $[1,2,3]$。全分辨率和 $1/2$ 尺度使用
$3\times3$ depthwise spatial mixing 与稠密 $1\times1$ gated channel mixer，便于 4090
Tensor Core 批量执行；$1/4$ 和 $1/8$ 尺度再并行混合局部 $3\times3$、
$1\times11+11\times1$ 轴向大核与空洞卷积。这一 clean-room 组合受 ConvIR 的
多尺度卷积恢复、OKNet 的 omni-kernel 及 DarkIR 的多尺度空间建模启发
\cite{convir,oknet,darkir}，但没有复制它们的预训练权重或实验数字。

\begin{figure}[H]
\centering
\resizebox{0.98\textwidth}{!}{%
\begin{tikzpicture}[
  node distance=7mm and 9mm,
  io/.style={draw,rounded corners,align=center,fill=CampGreen!12,text width=23mm,minimum height=12mm},
  route/.style={draw,rounded corners,align=center,fill=CampBlue!10,text width=38mm,minimum height=18mm},
  fusion/.style={draw,rounded corners,align=center,fill=CampPurple!14,text width=38mm,minimum height=15mm},
  head/.style={draw,rounded corners,align=center,fill=CampOrange!16,text width=38mm,minimum height=15mm},
  arr/.style={-{Latex[length=2.2mm]},thick}
]
\node[io] (x) {DAPI + Sobel\\$256\times256$};
\node[route,right=12mm of x,yshift=23mm] (naf) {Route A: NAF 主干\\四尺度局部/结构特征};
\node[route,right=12mm of x] (heavy) {Route B: Heavy Omni U\\四尺度解码细胞细节};
\node[route,right=12mm of x,yshift=-23mm] (freq) {Route C: Adaptive FFT\\学习 low/high 与双向调制};
\node[fusion,right=16mm of heavy] (fuse) {$1,1/2,1/4$: 双路 gate\\$1/8$: 三路 softmax + 2 SDPA};
\node[head,right=of fuse] (head) {ROI context + prototype\\decoder + adapter\\base/detail + deep supervision};
\node[io,right=of head] (y) {CD68\\prediction};
\draw[arr] (x)--(naf); \draw[arr] (x)--(heavy); \draw[arr] (x)--(freq);
\draw[arr] (naf)--(fuse); \draw[arr] (heavy)--(fuse); \draw[arr] (freq)--(fuse);
\draw[arr] (fuse)--(head)--(y);
\end{tikzpicture}}
\caption{v4 的三路主线。“并行”指三类特征在同一 autograd 计算图中分别计算，
再由内容相关门控融合；不是三个进程各占一块 GPU。}
\label{fig:v4-triroute}
\end{figure}

\paragraph{自适应频率路线。}
v3 对每张图都使用固定 cutoff $0.35$。v4 在这个稳定先验上根据当前 patch
预测有界的 cutoff 偏移和 transition 宽度，然后让池化的 high band 调制 low band，
low band 反过来调制 high band。这一“先分频、再双向交互”的思路受 AdaIR
\cite{adair} 启发，但要特别注意：AdaIR 处理噪声、雨、雾等退化，本项目是
DAPI$\rightarrow$CD68 跨 marker 映射，所以这只是候选归纳偏置，不能由论文结果推导一定涨分。
FFT 仍固定使用 float32，学习投影仍享受 BF16 AMP。v4 调用
\code{ParallelSpatialFrequencyMixer.forward_route()} 取得不含输入 identity 的 spatial--frequency update，
再将它作为真正独立的 Route C 送入三路融合；旧 v3/非三路配置仍调用兼容的
\code{forward(x)=x+forward_route(x)}，所以旧 checkpoint 行为和 state keys 不变。

\paragraph{内容感知融合。}
对于第 $r$ 条路的特征 $\mathbf F_r$，代码先做 LayerNorm2d 与 $1\times1$ 通道对齐，
再同时根据全局通道描述子和局部空间特征预测 logits：
\begin{equation}
 w_r(c,h,w)=\operatorname{softmax}_{r}\!\left(g_r^{\mathrm{channel}}(c)+
 g_r^{\mathrm{spatial}}(h,w)\right),\qquad \sum_r w_r=1.
\end{equation}
$1,1/2,1/4$ 尺度融合 NAF 与 Heavy U；$1/8$ 尺度融合 NAF、Heavy U 与
adaptive spatial-frequency。瓶颈后只使用两层原生 PyTorch
scaled-dot-product attention 做 Restormer 式转置通道注意力，避免在 $256^2$ token 上构建
昂贵的全局 Transformer。低分辨率 attention 是对全局交互的有界投资；ESC 对
高效 attention/卷积替代的分析进一步支持“不在高分辨率机械堆标准 Transformer”
的工程选择\cite{esc}。

融合不使用严格零初始化，而是以 $\alpha=0.05$ 小残差起步；频率内路尺度初值为 1，
稳定性由外层 $\alpha$ 统一控制。这避免“内层尺度等于 0，外层尺度也等于 0”导致分支
前几步完全收不到梯度。实际 BF16 forward/backward 检查已确认 Heavy U、adaptive frequency、
fusion 和 context 都有非零梯度。

\paragraph{日志如何证明分支真的参与。}
\code{RestorationOutput.fusion\_diagnostics} 每次前向返回四个尺度的 route gate 均值、
归一化熵、Heavy/Main 与 Frequency/Main 范数比、独立 frequency route/input RMS 比、融合尺度，以及 adaptive cutoff、
transition 和 low/high modulation gate。Trainer 不会每 batch 调用 \code{item()} 强制 GPU 同步；
它在设备上用 float32 累加，epoch 末统一写入
\path{log/<run>/epoch_metrics.jsonl} 和 \path{performance_summary.json}。这些值不是竞赛性能指标，
而是诊断“门控是否塌缩到某一路”、“新分支是否一直近似零”的工程证据。

\paragraph{为什么不物理删除 Restormer、CAMP-VS v2 和 ensemble。}
这些文件在 v4 YAML 下不会被实例化，因此不占用 v4 的 GPU 显存或每轮计算时间；
物理删除只会省下很小的源码空间，却会破坏旧 checkpoint 读取、配置回放、消融和回滚。
ensemble 更不是单模型 forward 中的一个开关，而是至少两个独立 checkpoint 训练完成后的
推理融合阶段。完整 Mamba 本轮也不进主线：它会引入 selective-scan CUDA 扩展、环境风险，
而“理论 FLOPs 较低”不自动等于 4090 实际吞吐更高。

\subsection{输入预处理与 Sobel 结构提示}

当前配置启用 \code{use_sobel_input=true}。固定 Sobel 卷积核计算水平和垂直梯度，再取

\begin{equation}
  S=\sqrt{G_x^2+G_y^2+\varepsilon},
\end{equation}

得到边缘强度图。当前 Dataset 已把 RGB 存储文件转换为一个逻辑灰度通道，所以 Sobel 直接在这一通道上
计算。原输入和 Sobel 图在通道维拼接成 2 通道，再通过保持 $256\times256$ 尺寸的 $3\times3$ stem 卷积。Sobel
没有可训练参数，它只是明确告诉模型“哪里变化快”，类似在原图旁边附上一张铅笔勾边图。

\subsection{NAF 局部编码器}

\term{NAF}{Nonlinear Activation Free}来自 NAFNet 图像恢复思想\cite{nafnet}。名称的重点是：
NAFBlock 主干不依赖 ReLU/GELU 这类常规逐点激活，而使用乘法门控建立非线性。它并不意味着整个
模型完全没有激活函数；当前 task adapter 使用 GELU，prototype 分配使用 softmax，最终输出使用 sigmoid。

\subsubsection{SimpleGate}

若输入特征有 $2C$ 个通道，SimpleGate 沿通道一分为二：

\begin{equation}
  \operatorname{SimpleGate}([\mathbf{x}_1,\mathbf{x}_2])
  =\mathbf{x}_1\odot\mathbf{x}_2,
  \qquad \mathbf{x}_1,\mathbf{x}_2\in\mathbb{R}^{C\times H\times W}.
\end{equation}

$\odot$ 表示逐元素相乘。直观地说，一半特征像“内容”，另一半像“开关强度”，两者相乘后共同决定
信息是否通过。它没有额外参数，但乘法本身产生了非线性关系。

\subsubsection{SCA 简化通道注意力}

\term{SCA}{Simplified Channel Attention}先对每个通道做全局平均池化，得到 $C\times1\times1$ 的摘要；
再用 $1\times1$ 卷积生成通道缩放系数，并乘回原特征：

\begin{equation}
  \operatorname{SCA}(\mathbf{x})=\mathbf{x}\odot
  \operatorname{Conv}_{1\times1}(\operatorname{GAP}(\mathbf{x})).
\end{equation}

它比包含多层全连接和 sigmoid 的传统 SE 模块更简单。这里的“注意力”可理解为模型在问：当前图像中，
哪些特征通道更值得放大？

\subsubsection{NAFBlock 的真实顺序}

图~\ref{fig:naf-block}按 \code{naf_blocks.py} 展示两个残差子块。第一支负责空间混合与通道重标定，
第二支像门控前馈层。$\beta$ 与 $\gamma$ 都从 0 初始化，所以网络刚建立时每个 block 接近恒等映射，
有利于深层模型稳定开始训练。

\begin{figure}[H]
\centering
\resizebox{0.95\textwidth}{!}{%
\begin{tikzpicture}[
  node distance=5mm,
  op/.style={draw,rounded corners,fill=CampBlue!9,align=center,minimum height=10mm,text width=22mm},
  gate/.style={draw,rounded corners,fill=CampOrange!14,align=center,minimum height=10mm,text width=21mm},
  add/.style={draw,circle,fill=CampGreen!14,minimum size=8mm,inner sep=0pt},
  arr/.style={-{Latex[length=2mm]},thick}
]
\node[op] (x) {输入 $x$};
\node[op,right=of x] (ln1) {LayerNorm\\2d};
\node[op,right=of ln1] (c1) {$1\times1$ Conv\\通道扩张};
\node[op,right=of c1] (dw) {$3\times3$\\Depthwise Conv};
\node[gate,right=of dw] (g1) {SimpleGate};
\node[gate,right=of g1] (sca) {GAP+$1\times1$\\SCA};
\node[op,right=of sca] (p1) {$1\times1$ Conv\\$\times\beta$};
\node[add,right=of p1] (a1) {$+$};
\draw[arr] (x)--(ln1)--(c1)--(dw)--(g1)--(sca)--(p1)--(a1);
\draw[arr] (x) to[bend right=28] (a1);

\node[op,below=13mm of ln1] (ln2) {LayerNorm\\2d};
\node[op,right=of ln2] (ff) {$1\times1$ Conv\\通道扩张};
\node[gate,right=of ff] (g2) {SimpleGate};
\node[op,right=of g2] (p2) {$1\times1$ Conv\\$\times\gamma$};
\node[add,right=of p2] (a2) {$+$};
\node[op,right=of a2] (y) {输出 $y$};
\draw[arr] (a1) |- (ln2);
\draw[arr] (ln2)--(ff)--(g2)--(p2)--(a2)--(y);
\draw[arr] (a1) to[bend left=24] (a2);
\end{tikzpicture}}
\caption{NAFBlock：两次归一化、两条门控残差路径和零初始化残差缩放}
\label{fig:naf-block}
\end{figure}

\subsubsection{四个编码尺度}

\code{MultiMarkerRestorer._encode()} 的 stem 保持 $256\times256$，每个 NAF stage 之后（最后一级除外）
再用 $2\times2$、stride 2 卷积下采样。当前通道为 $[64,128,256,512]$，depths 为 $[2,2,4,6]$：

\begin{table}[H]
\centering
\caption{局部编码器各尺度}
\begin{tabular}{ccccc}
\toprule
相对尺度 & 空间尺寸 & 通道数 & NAFBlock 数 & 主要作用 \\
\midrule
1 & $256\times256$ & 64 & 2 & 细边缘与局部纹理 \\
1/2 & $128\times128$ & 128 & 2 & 中等范围结构 \\
1/4 & $64\times64$ & 256 & 4 & 更大感受野与 skip \\
1/8 & $32\times32$ & 512 & 6 & 原型混合前的瓶颈表示 \\
\bottomrule
\end{tabular}
\end{table}

\subsection{瓶颈处的共享/任务余弦原型混合}

原型混合器本身没有被 Restormer 替换；它仍调用 \code{MultiTaskPrototypeMixer}。v4 只在三路特征进入
原型层之前增加两层低分辨率原生 PyTorch SDPA 精炼，因此全局交互和 shared/task prototype 是先后互补而非
两套互斥主模型。随后把
$\mathbf{F}\in\mathbb{R}^{B\times512\times32\times32}$ 展平后，每个空间位置成为一个 512 维 token，
总计 1024 个 token。代码把 token、8 个 shared prototype 和 8 个 CD68 task prototype 都做 $L_2$
归一化，再计算余弦相似度和温度为 0.1 的 softmax：

\begin{equation}
  A_{ik}^{(s)}=\operatorname{softmax}_k
  \left(\frac{\bar{\mathbf f}_i^{\mathsf T}\bar{\mathbf p}_k^{(s)}}{0.1}\right),
  \qquad
  A_{ik}^{(t)}=\operatorname{softmax}_k
  \left(\frac{\bar{\mathbf f}_i^{\mathsf T}\bar{\mathbf p}_k^{(t)}}{0.1}\right).
\end{equation}

shared 与 task attention 分别对原型加权求和；两个 context 拼接后经线性层投影回 512 维，再以可学习
\code{fusion_scale} 残差加回原特征。该尺度初始值为 0.1，所以原型分支从小幅修正开始，而不是一开始
覆盖卷积特征。相似度、归一化和 softmax 在 float32 中计算，随后再转回输入 dtype。

\begin{figure}[H]
\centering
\resizebox{0.96\textwidth}{!}{%
\begin{tikzpicture}[
  node distance=6mm,
  op/.style={draw,rounded corners,fill=CampPurple!10,align=center,minimum height=11mm,text width=24mm},
  add/.style={draw,circle,fill=CampGreen!14,minimum size=8mm,inner sep=0pt},
  arr/.style={-{Latex[length=2mm]},thick}
]
\node[op] (x) {$32\times32$ 特征\\展平并 $L_2$ 归一化};
\node[op,right=of x] (shared) {8 shared\\prototypes};
\node[op,below=11mm of shared] (task) {8 CD68 task\\prototypes};
\node[op,right=of shared] (attn1) {cosine / 0.1\\shared softmax};
\node[op,right=of task] (attn2) {cosine / 0.1\\task softmax};
\node[op,right=11mm of attn1] (concat) {加权 context\\拼接并线性投影};
\node[op,right=of concat] (scale) {$\times$ learnable\\fusion scale};
\node[add,right=of scale] (add) {$+$};
\draw[arr] (x)--(shared)--(attn1)--(concat)--(scale)--(add);
\draw[arr] (x) |- (task); \draw[arr] (task)--(attn2)--(concat);
\draw[arr] (x.south) -- ++(0,-22mm) -| (add.south);
\end{tikzpicture}}
\caption{MultiTaskPrototypeMixer：共享/任务原型注意力与残差融合}
\label{fig:prototype-mixer}
\end{figure}

\subsection{解码器与 skip connection}

解码器有三个 \code{DecoderStage}。每一级先用双线性插值把低分辨率特征放大到 skip 的尺寸，再用
$1\times1$ 卷积调整通道，把放大特征和编码器 skip 在通道维拼接，经过另一个 $1\times1$ 卷积融合，
最后用 NAFStage 精炼。skip connection 的作用像给“根据摘要重画原图”的人补充原始草稿：瓶颈保留
大局，浅层 skip 带回精细边缘。

三个解码输出尺度分别是 $64\times64$、$128\times128$ 和 $256\times256$；每级 NAFStage 都有
2 个 NAFBlock。每一级随后可经 task adapter，并由独立 $1\times1$ head 产生该尺度的 sigmoid 预测。

\subsection{marker conditioning、FiLM 与 adapter}

\term{conditioning}{条件控制}让共享模型知道当前要预测哪个 marker。当前主模型为目标建立 64 维
task embedding；它不调用 CAMP-VS v2 的 organ conditioner。
\term{FiLM}{Feature-wise Linear Modulation}根据 embedding 生成逐通道缩放和偏置：

\begin{equation}
  \operatorname{FiLM}(\mathbf{F})=(1+\boldsymbol{\gamma})\odot\mathbf{F}
  +\boldsymbol{\beta}.
\end{equation}

解码器每级的 \code{TaskAdapter} 先由 embedding 生成 scale/shift，线性层从零初始化；再经过
LayerNorm2d、depthwise $3\times3$、GELU 和 $1\times1$ pointwise 路径，残差尺度也从 0 开始。
当前只有 CD68 一个任务，embedding 不能在多个 marker 之间做选择，但接口保留了以后多任务扩展能力。

\subsection{direct output 与零初始化 base/detail residual}

v3/v4 都保留每个解码尺度的 direct $1\times1$ output head；当前 v4 在最终 $256\times256$ 尺度开启
\code{LaplacianBaseDetailHead} 的 residual-to-reference 模式。direct 路径仍为

\begin{equation}
  \widehat{\mathbf y}_s=\sigma\!\left(\operatorname{Conv}_{1\times1}(\mathbf F_s)\right),
  \qquad s\in\{64,128,256\}.
\end{equation}

最终尺度再按上一节公式加入零初始化 coarse/detail residual。全局强度标定器
\path{GlobalIntensityCalibrator} 仍未启用，因为本轮没有验证证据表明它能稳定提高 JPG 指标。

\subsection{Deep Supervision}

\term{深监督}{Deep Supervision}是在中间解码层也产生预测，让浅层和深层都直接收到训练信号。
当前只有三个输出尺度：$256$、$128$、$64$。目标图用 area 插值缩放到对应大小。损失权重在训练中
从 $[1.0,0.5,0.25]$ 平滑退火到 $[1.0,0.1,0.05]$：前期让中间层容易获得梯度，后期把容量集中到
真正用于提交的全分辨率输出。

\subsection{当前实际模型规模}

结构统计区分稳定 v2 对照、已停止的 v3 和当前 v4 候选。v2 为 16,232,644 参数、
约 39.160 GMAC；v3 为 18,074,199 参数、约 44.11 GMAC；v4 的真实构造结果为
29,130,894 参数、约 75.997 GMAC/单张 $256\times256$ forward。
MAC 只统计可计数卷积/线性运算，不完整包含 FFT、SDPA attention matmul、SSIM、JPEG 解码、15\% batch 的第二次 D4 forward/backward
或九张邻域图像的 I/O，因此不能直接换算成小时数。

\begin{table}[H]
\centering
\caption{稳定对照、v3 与 v4 候选的结构规模}
\begin{tabular}{lrrr}
\toprule
配置 & 参数量 & 约 GMAC/图 & 状态 \\
\midrule
retrain-v2 稳定对照 & 16,232,644 & 39.160 & 已完成训练/验证 \\
max-v3 已停止候选 & 18,074,199 & 44.11 & 用于结构对照 \\
其中轻量 Detail U & 72,016 & 0.589 & v3 已启用 \\
max-v4 当前候选 & 29,130,894 & 75.997 & 工程 smoke 通过，待 4090 训练 \\
其中 Heavy Omni U & 4,509,672 & 包含于 v4 & v4 已启用 \\
多尺度+三路融合 & 6,453,341 & 包含于 v4 & v4 已启用 \\
\bottomrule
\end{tabular}
\label{tab:model-params}
\end{table}

\important{v4 当前\textbf{启用} Sobel、NAF、Heavy Omni U、adaptive spatial/frequency、四尺度 gate、
$1/8$ tri-path SDPA、verified $3\times3$ ROI context、9-token cross-attention、shared/task prototypes、CD68 adapters 和 residual base/detail。
故意关闭的模块包括与单目标不相容的 FAMO/correlation、与零配准偏移冲突的 shift-tolerant loss、
没有正收益证据的 foreground loss，以及会显著扩大时间/依赖风险的完整 Mamba、扩散和 GAN。这里的
“全开”指所有与单目标、已对齐、9 小时预算相容的候选模块，而不是机械打开互相冲突的开关。}

\section{数据流向详解}

\subsection{用一个真实 patch 对照源码路径}

图~\ref{fig:real-code-pipeline}继续使用 \code{ROI007_09_11}。上方蓝色箭头表示 DAPI 真正进入模型的
路径；下方红色分支表示训练时才会读取的 CD68 真值。图中的 CD68 小图是\textbf{训练标签，不是模型
预测}。当前报告构建目录没有完成正式训练的 competition checkpoint，因此本文不会拿真值冒充“生成结果”。
真正的预测外观必须由指定 checkpoint 在 test DAPI 上运行后产生。

\begin{figure}[H]
\centering
\includegraphics[width=\textwidth]{docs/essay/figures/real_patch_code_pipeline.png}
\caption{一个真实配对 patch 从 JPEG、Dataset、MultiMarkerRestorer 到损失与提交文件的代码路径。}
\label{fig:real-code-pipeline}
\end{figure}

表~\ref{tab:real-patch-code-path}把图中的每个方框与当前源码对应起来。需要特别注意：四个 target 文件夹
不会同时“喂给”单目标 CD68 模型。配置中的 \code{data.targets} 决定 Dataset 读取哪些真值，
\code{model.name} 决定 CLI 构造哪一个网络；当前初赛配置只读取 CD68 作为监督目标，并构造一个
\code{MultiMarkerRestorer} 实例。

\begin{table}[H]
\centering
\caption{真实 patch 在代码中的表示变化}
\label{tab:real-patch-code-path}
\small
\begin{tabularx}{\textwidth}{>{\raggedright\arraybackslash}p{0.13\textwidth}
                              >{\raggedright\arraybackslash}p{0.17\textwidth}
                              >{\raggedright\arraybackslash}p{0.29\textwidth}X}
\toprule
阶段 & 输入/数据键 & 对应源码 & 产生的结果 \\
\midrule
图像解码 & DAPI/CD68 JPG & \path{utils/image_io.py} 的 \path{load_image_tensor()} & RGB uint8 HWC 的三个相同通道先取均值并四舍五入，变为 $[0,1]$ 的单通道 float32 CHW。\\
样本组装 & manifest 一行 & \path{data/dataset.py} 的 \path{VirtualStainingDataset.__getitem__()} & 返回 \path{batch["input"]}、\path{batch["targets"]["CD68"]} 和 ROI/stem 元数据。\\
选择模型 & 合并后的 YAML & \path{cli.py} 的 \path{_build_model()} & 当 \path{model.name=multi_marker_restorer} 时实例化当前 v4 \path{MultiMarkerRestorer}；不会另外实例化 CAMP-VS v2 或独立 Residual U-Net。\\
前向传播 & $B\times1\times256\times256$ DAPI 与 9 个 context tiles & \path{models/multi_marker_restorer.py} 的 \path{MultiMarkerRestorer.forward()} & NAF、Heavy U、自适应 FFT 分别提取特征，经四尺度门控、低分辨率 SDPA、context、原型与 task-adapted 解码后返回 \path{RestorationOutput.predictions["CD68"]}。\\
训练比较 & 预测图与 CD68 真值 & \path{losses/scheduled_composite.py}、\path{losses/competition_proxy.py} 和 \path{engine/trainer.py} & 两阶段重建/结构损失、深监督退火、后段 float metric proxy 与可选 D4 等变约束；CD68 真值只走监督分支，不作为模型输入。\\
推理保存 & test DAPI 与 checkpoint & \path{engine/inferencer.py}、\path{submission/writer.py} & 预测先保存为 \path{<stem>.jpg}，提交构建时命名为 \path{<stem>_fake.jpg} 并写入 ZIP。\\
\bottomrule
\end{tabularx}
\end{table}

从神经网络内部再向下看，\code{MultiMarkerRestorer.forward()} 不是直接对每个像素套一个公式。它先调用
\code{_encode()} 得到四尺度 features，在瓶颈执行 shared/task prototype mixing，再按任务调用
\code{_decode_task()}。CD68 分支产生主预测、三尺度深监督和 prototype attention 诊断张量，最后装入
\code{RestorationOutput}。训练器取预测与真值算 loss；推理器只取预测，
执行可选 D4 平均、截断到 $[0,1]$、四舍五入到 uint8，并以 quality 100、subsampling 0 编码一次 JPEG。
提交 writer 随后直接复制已有 JPEG 并改成比赛文件名，避免第二次有损编码。

\subsection{训练数据流}

图~\ref{fig:train-flow}把训练过程从硬盘一直画到参数更新。张量形状按当前单逻辑通道配置表示。

\begin{figure}[H]
\centering
\resizebox{\textwidth}{!}{%
\begin{tikzpicture}[
  node distance=6mm,
  cpu/.style={draw,rounded corners,fill=CampGreen!11,align=center,minimum height=13mm,text width=25mm},
  gpu/.style={draw,rounded corners,fill=CampBlue!11,align=center,minimum height=13mm,text width=27mm},
  loss/.style={draw,rounded corners,fill=CampOrange!15,align=center,minimum height=13mm,text width=29mm},
  arr/.style={-{Latex[length=2mm]},thick}
]
\node[cpu] (jpg) {DAPI/CD68 JPG\\manifest 路径};
\node[cpu,right=of jpg] (workers) {8 workers 并发\\PIL 解码};
\node[cpu,right=of workers] (tensor) {1 通道 CHW float32\\归一化到 $[0,1]$};
\node[cpu,right=of tensor] (aug) {成对几何增强\\DAPI 强度增强};
\node[cpu,right=of aug] (pin) {batch + pin memory\\{\scriptsize $[12,1,256,256]$}};
\draw[arr] (jpg)--(workers)--(tensor)--(aug)--(pin);

\node[gpu,below right=15mm and -15mm of jpg] (copy) {non-blocking H2D\\channels\_last};
\node[gpu,right=of copy] (model) {MultiMarker\\Restorer forward};
\node[loss,right=of model] (criterion) {调度复合损失\\主输出+深监督};
\node[gpu,right=of criterion] (back) {autograd backward\\可选 D4 一致性};
\node[gpu,right=of back] (step) {AdamW step\\clip + EMA};
\draw[arr] (pin)--(copy); \draw[arr] (copy)--(model)--(criterion)--(back)--(step);
\draw[arr,dashed] (step) to[bend right=28] (model);
\end{tikzpicture}}
\caption{训练时从 JPG 到模型参数更新的数据流}
\label{fig:train-flow}
\end{figure}

进一步拆解如下：

\begin{enumerate}
  \item AutoDL v4 默认由 8 个 persistent workers 按 manifest 预取配对 JPG，PIL 解码为 RGB uint8；Windows 可显式覆盖为 0。
  \item \code{load_image_tensor()} 对三个相同 RGB 通道取均值并四舍五入为一个逻辑通道，再转为 CHW float32 并除以 255。
  \item PairedTransform 同步改变 DAPI/CD68 的几何，仅扰动 DAPI 强度。
  \item DataLoader 把 6 个样本堆成物理 batch；每累积 2 个 batch 才更新一次参数，所以 effective batch 仍为 12。页锁定内存为 CPU 到 GPU 传输做准备。
  \item \code{move_to_device()} 对张量使用 non-blocking CUDA copy；训练器在启用时把 4D 浮点输入转为 channels-last 物理布局。
  \item autocast 选择 BF16 或 FP16 执行适合低精度的算子；模型返回结构化预测结果。
  \item scheduled composite loss 将预测与 CD68 真值比较，连续调度均方误差、稳健像素误差和 SSIM 等权重，并把 256/128/64 三个分辨率的退火深监督合并；后段再加入小权重 float metric proxy。
  \item 正常每个 DataLoader batch 执行一次反向传播、梯度累积、优化器更新和 EMA 更新；v4 的 15\% batch 另计算一次随机 D4 等变一致性。CUDA OOM 时只在 batch 内切 chunk，不改变有效 batch。
\end{enumerate}

\subsection{验证数据流}

验证不更新模型。它要回答的问题是：“用当前权重预测从未用于梯度更新的 ROI，文件质量怎么样？”

\begin{figure}[H]
\centering
\resizebox{0.98\textwidth}{!}{%
\begin{tikzpicture}[
  node distance=7mm,
  box/.style={draw,rounded corners,fill=CampGray,align=center,minimum height=12mm,text width=27mm},
  metric/.style={draw,rounded corners,fill=CampPurple!12,align=center,minimum height=12mm,text width=28mm},
  arr/.style={-{Latex[length=2mm]},thick}
]
\node[box] (valjpg) {val DAPI/CD68\\无随机增强};
\node[box,right=of valjpg] (eval) {model.eval\\inference mode};
\node[box,right=of eval] (pred) {预测 clamp\\到 $[0,1]$};
\node[metric,above right=8mm and 8mm of pred] (float) {float 域\\SSIM/PSNR};
\node[metric,right=of pred] (u8) {uint8 域\\round/clip};
\node[metric,below right=8mm and 8mm of pred] (jpgm) {JPG 域\\预测编码再解码};
\node[box,right=18mm of u8] (record) {逐图记录\\stem/ROI/分层};
\node[box,right=of record] (agg) {per-target + macro\\ROI 聚合};
\draw[arr] (valjpg)--(eval)--(pred);
\draw[arr] (pred)--(float); \draw[arr] (pred)--(u8); \draw[arr] (pred)--(jpgm);
\draw[arr] (float)--(record); \draw[arr] (u8)--(record); \draw[arr] (jpgm)--(record);
\draw[arr] (record)--(agg);
\end{tikzpicture}}
\caption{验证时的三域指标与 ROI 聚合数据流}
\label{fig:validation-flow}
\end{figure}

验证 loader 不创建 PairedTransform；模型使用 \code{eval()} 和 \code{torch.inference_mode()}，所以 dropout
等模块进入推理行为，也不保存梯度图。每张预测先截断到 $[0,1]$，再进入三个指标域。结束后恢复模型
原来的 train/eval 状态，避免验证永久改变训练模式。

\subsection{测试与推理数据流}

测试 manifest 只有 DAPI。Inferencer 保证 loader 不能随机打乱，然后按稳定顺序预测。当前正式提交配置为
\code{tta=d4}、\code{tta\_policy=configured}，因此配置中的 D4 是权威选择，不会被一次不完整的自动晋级
审计静默关闭。同一 batch 会经历恒等、$90^\circ$、$180^\circ$、$270^\circ$、水平翻转、垂直翻转、转置和反转置八种
变换。每个预测经过精确逆变换后在 float32 中求平均：

\begin{equation}
  \widehat{\mathbf y}_{\mathrm{D4}}
  =\frac{1}{8}\sum_{k=1}^{8}T_k^{-1}\left(f_\theta(T_k(\mathbf x))\right).
\end{equation}

平均结果 round 到 uint8，并按 checkpoint ImageSpec 指定的 RGB 模式，以 quality 100、subsampling 0、
optimize false 编码一次 JPEG，先保存为 \code{predictions/CD68/<stem>.jpg}。提交构建器再把它直接复制成
\code{results/test/CD68/<stem>_fake.jpg}，不会把 JPEG 再解码重编码。

推理也有最多三次定向 CUDA OOM 降级：发生明确 CUDA out-of-memory 时把 batch 切成更小 chunk；其他
异常不会被误当 OOM 吞掉。

\section{训练详解}

\subsection{epoch 与 batch：把“看很多遍数据”说清楚}

\term{epoch}{轮次}表示把当前训练集完整遍历一次。当前有 5004 个训练 patch，v4 物理 batch size 为 6，
DataLoader 不丢弃最后不足 6 张的 batch，因此每个 epoch 大约有

\begin{equation}
  \left\lceil\frac{5004}{6}\right\rceil=834
\end{equation}

个 batch。当前配置最多训练 100 epoch，不代表“只看 100 张图”，而是反复遍历数据 100 次，并且每次都可能
得到新的随机增强。

\code{Trainer.fit()} 是外层 epoch 循环，\code{train_epoch()} 是内层 batch 循环。每个 epoch 开始时模型
进入 train 模式、清空梯度、重置显存峰值统计；结束时记录真实耗时、吞吐量、峰值显存、RSS 内存、学习率、
OOM 次数和损失分量。

\subsection{一次 batch 内发生什么}

对第 $t$ 个 batch，训练器完成：

\begin{enumerate}
  \item \code{unpack_batch()} 得到输入、目标字典和元数据；
  \item 张量移动到 GPU，并按配置转 channels-last；
  \item 在 autocast 中执行模型 forward 与 loss；
  \item 检查总 loss 必须是一个有限标量，防止 NaN/Inf 污染参数；
  \item 将 loss 除以梯度累积因子并 backward；
  \item 到达累积边界后进行梯度裁剪、AdamW step、EMA 更新，并以 set-to-none 方式清空梯度。
\end{enumerate}

\subsection{梯度累积与 effective batch}

当前 \code{batch_size=6}、\code{gradient_accumulation=2}，所以正常情况下

\begin{equation}
  B_{\mathrm{effective}}=6\times2=12.
\end{equation}

每个 epoch 正常约有 834 个 DataLoader batch、417 次 optimizer step，EMA 也大约更新 417 次。若 24 GiB 4090 发生明确 CUDA OOM，
训练器会把同一个 DataLoader batch 内部分成 3+3 等 chunk；各 chunk 的梯度按真实样本数占原 batch 的比例缩放后再做一次
optimizer step，所以有效 batch 仍是 12，不会因为降级偷偷变成 24。

\subsection{AMP 混合精度}

\term{AMP}{Automatic Mixed Precision，自动混合精度}让卷积和矩阵乘等适合的算子使用 BF16/FP16，
对需要稳定范围的算子保留 float32。这样通常可以减少显存并利用 Tensor Core 加速\cite{pytorch_amp}。

当前 v4 显式设置 \code{amp=true}、\code{amp_dtype=bfloat16}，面向 RTX 4090/5090 使用 BF16 Tensor Core。
GradScaler 只在 CUDA+FP16 时启用；BF16 的指数范围较大，代码不启用 scaler。SSIM、FFT、
Sobel gradient 和 prototype cosine 等敏感计算会显式升到 float32，这说明“开 AMP”不是把每个数都粗暴降精度。

\subsection{AdamW、学习率、warmup 与梯度裁剪}

当前优化器是 AdamW，初始学习率 $1.8\times10^{-4}$，weight decay $10^{-4}$。AdamW 将 weight decay
与自适应梯度更新解耦\cite{adamw,pytorch_adamw}。前 4 个 epoch 是 warmup，学习率从较小倍数逐步升到
基础值；之后使用余弦函数下降。v4 还把 bias、normalization、NAF beta/gamma 和 residual scale 放入
no-decay 参数组；CUDA 支持时使用 fused AdamW，不支持则自动回退普通 AdamW：

\begin{equation}
  \eta(p)=\eta_0\frac{1+\cos(\pi p)}{2},\qquad p\in[0,1].
\end{equation}

前 4 个 epoch 的 warmup 像汽车起步时缓慢踩油门，避免随机初始化阶段更新过猛。每次 optimizer step 前，所有优化参数的
梯度总范数被裁剪到 1.0，防止个别 batch 产生异常大的更新。

\subsection{EMA：不是另一个独立模型}

\term{EMA}{Exponential Moving Average，指数移动平均}保存一份参数影子：

\begin{equation}
  \bar\theta_t=0.999\bar\theta_{t-1}+0.001\theta_t.
\end{equation}

它不是再训练一遍模型，而是在每个 optimizer step 后平滑跟踪当前参数。验证 EMA 时，代码暂时备份 raw
权重、把 EMA shadow 复制进模型、完成验证后恢复 raw。短训练时 EMA 可能严重滞后，所以当前训练配置
明确在每个验证 epoch 同时评价 \code{raw} 和 \code{ema}，按 JPG SSIM、再按 JPG PSNR 选择更好的来源，
而不是先验认定 EMA 一定更优。选中来源会写入 checkpoint 的 extra 元数据，提交时
\code{weight_source=best_jpg} 会读取它。

\subsection{稳定 v2 对照的固定权重复合损失}

只用像素误差可能得到模糊结果，只用结构指标又可能忽略强度与数值误差。稳定 v2 对照的
\code{loss.schedule=constant}，所以 \code{_build_loss()} 构造 \code{CompositeRestorationLoss}，
不会调用 two-phase 的 \code{ScheduledCompositeLoss}。每个输出尺度组合五项重建损失：

\begin{align}
  \mathcal L_{\mathrm{Charb}} &= \frac{1}{N}\sum_i w_i
  \sqrt{(\hat y_i-y_i)^2+\varepsilon^2},\\
  \mathcal L_{\mathrm{SSIM}} &= 1-\operatorname{SSIM}(\hat{\mathbf y},\mathbf y),\\
  \mathcal L_{\mathrm{freq}} &= \left\|\log(1+|\mathcal F(\hat{\mathbf y})|)
  -\log(1+|\mathcal F(\mathbf y)|)\right\|_1.
\end{align}

其中 $w_i$ 是由目标局部均值差异构造的 structure weight；另外还有 MS-SSIM 和 Sobel gradient。
Charbonnier 是平滑的 L1 近似，gradient 强调边缘，frequency 比较 FFT 的 log amplitude 而不直接约束
容易不稳定的相位。

\begin{table}[H]
\centering
\caption{稳定 v2 对照的固定损失重建项权重}
\begin{tabular}{lrp{0.54\textwidth}}
\toprule
分量 & 权重 & 直观作用 \\
\midrule
Charbonnier & 0.40 & 带 structure weight 的稳健像素重建 \\
SSIM & 0.35 & 局部结构、亮度与对比度 \\
MS-SSIM & 0.10 & 多尺度结构 \\
Gradient & 0.10 & Sobel 边缘一致 \\
Frequency & 0.05 & 频谱幅度和纹理能量一致 \\
\bottomrule
\end{tabular}
\label{tab:loss-weights}
\end{table}

三个 deep-supervision 尺度各自计算上述五项，再按 $[1.0,0.5,0.25]$ 相加。MSE、pyramid 与 statistics
权重在有效配置中均为 0。跨 marker correlation 的配置权重虽为 0.02，但当前只有 CD68 一个 prediction，
\code{task_correlation_loss()} 按定义返回图连接的 0，不会贡献梯度。

prototype activation 与 diversity 两项则\textbf{真实启用}，权重各为 0.001。前者让采样到的瓶颈 token
靠近至少一个 shared/CD68 prototype，后者惩罚同一 bank 内原型间的平方余弦相似度，降低所有原型坍缩到
同一方向的风险。因此总损失可以概括为

\begin{equation}
  \mathcal L=\sum_{s\in\{256,128,64\}}\lambda_s\mathcal L_{\mathrm{recon},s}
  +0.001\mathcal L_{\mathrm{proto\text{-}act}}
  +0.001\mathcal L_{\mathrm{proto\text{-}div}}.
\end{equation}

\subsection{v3/v4 两阶段损失、指标对齐与 D4 等变}

v3 引入、v4 继承 \code{ScheduledCompositeLoss}。训练进度 $p\in[0,1]$ 从 phase A 平滑过渡到 phase B：前者
以 Charbonnier 和多尺度结构为主，后者提高 MSE/SSIM 对最终像素指标的约束，并以小权重加入
Laplacian pyramid、gradient、statistics 与 log-frequency。连续 cosine 插值避免在阶段边界突然改变
梯度方向。三个深监督权重也按同一进度从 $[1,0.5,0.25]$ 退火到 $[1,0.1,0.05]$。

从训练后 55\% 开始，最终 $256\times256$ 输出还加入权重上限 0.06 的 float-domain metric proxy。
它用与 scikit-image float SSIM 一致的 uniform $7\times7$、sample covariance 和 valid crop，并按每图
MSE 计算 capped PSNR loss。它\textbf{不模拟} uint8 round 或 JPEG 编解码；真正 checkpoint 选择仍以
非可微的 JPG round-trip 验证为准，不能把这项 loss 称为“可微 JPEG”。

训练进度超过 10\% 后，D4 consistency 权重平滑升至 0.015；v4 在 15\% 的 batch 随机抽取一个非恒等 D4 变换，
同步旋转/翻转中心图、九张 context tiles、valid mask 和二维 offsets，再最小化

\begin{equation}
  \mathcal L_{\mathrm{eq}}=
  \operatorname{SmoothL1}\left(T^{-1}f(Tx),\operatorname{stopgrad}(f(x))\right).
\end{equation}

这会平均增加约 15\% batch 的额外 forward/backward，但不增加推理参数。末段 hard-mix sampler 只在最后
22\% epoch 把 30\% 槽位温和分给最高 activity 桶；总采样数不变，test 永远不参与训练。

\subsection{channels-last、cuDNN benchmark 与 float32 matmul precision}

当前配置为 CUDA 训练启用 channels-last。逻辑张量形状仍写成 NCHW，但物理内存更接近 NHWC 排列，
卷积密集模型可能更容易走 Tensor Core 友好的内核。训练模型和输入都会转换；Validator 也转换输入。
当前 Inferencer 没有对应的 channels-last 输入转换路径，因此不能把训练/验证优化自动外推到最终推理。

\code{deterministic=false}、\code{cudnn_benchmark=true} 允许 cuDNN 针对固定 $256\times256$ 尺寸测试并
选择较快卷积算法，代价是严格逐位复现能力下降。\path{float32_matmul_precision=high} 在支持的 NVIDIA
GPU 上允许 float32 矩阵乘采用更快的内部精度路径（通常包括 TF32）。这些是性能设置，不是保证每台
GPU 都获得相同比例加速的魔法开关。

\subsection{进度条、异步指标、profiler 与 torch.compile}

\begin{itemize}
  \item \code{progress_bar=auto}：仅在交互式终端的 stderr 是 TTY 时显示 tqdm batch 进度，内容包括 loss、图/秒、学习率和峰值显存；nohup 重定向时默认关闭，避免日志被刷新字符污染。
  \item \code{async_metric_logging=true}：loss 分量留在 GPU 累积，epoch 末统一同步，避免每个 batch 因转 Python float 让 GPU 等待 CPU。
  \item 验证按 pass 输出静态进度：训练期 raw/EMA 分别标为 \code{pass=1/2}、\code{pass=2/2}；训练后的普通与 D4 验证也会打印开始、完成、样本数和耗时，不会在长时间 JPG round-trip 中完全沉默。
  \item \code{gpu_monitor.enabled=true}：后台每 2 秒调用一次 \code{nvidia-smi}，把 GPU utilization、显存、功耗和温度的 mean/p95/max 汇总进 epoch 日志；监控失败只丢诊断，不中断训练。
  \item \code{profiler.enabled=false}：可选 profiler 能跟踪最初若干 batch 并导出 Chrome trace，但当前生产训练不承担这项开销。
  \item \code{compile.enabled=false}：可选 \code{torch.compile} 失败会回退 eager；当前配置没有启用，也不存在“已经通过 compile 加速”的结果。
\end{itemize}

\subsection{OOM 恢复}

只有错误明确属于 CUDA out-of-memory 时才进入恢复：清梯度、清缓存，把同一个 DataLoader batch 在内部
切成更小 chunk。\textbf{配置的 gradient accumulation 不再倍增}；若 OOM 清除了一个尚未完成的累积窗口，
窗口计数会从重试成功的 batch 重新开始。最多三次；chunk 已经为 1 仍 OOM 时抛出错误。这样能保持原有效
batch，又不会把文件损坏、shape 错误等普通 bug 伪装成显存问题。虽然 YAML 有
\code{max_oom_retries=3}，当前 Trainer 实现实际把“三次”直接写在判断逻辑中。

\subsection{验证频率、checkpoint 与早停}

v4 有效 \code{validate_every_n_epochs=5}，所以每 5 个 epoch 验证一次；最后一个
epoch 仍会强制验证；跳过验证的 epoch 只更新 \code{last.ckpt}，不会冒充 best 或 top-k。

每个 epoch 原子写入 \code{last.ckpt}。验证指标创新高时分别写
\code{best_ssim.ckpt} 与 \code{best_proxy.ckpt}。当前 \code{save_top_k=0}，不额外保留重复 top-k 大文件。
checkpoint 版本 2 包含模型、EMA、优化器、scheduler、scaler、loss 状态、epoch、global step、完整配置、
manifest hash、ImageSpec、目标列表、指标历史、随机数状态和 DataLoader generator 状态，因此 resume
不是“只加载一份权重重新开始”。

early-stopping patience 为 5，计数单位是\textbf{实际验证事件}；即连续 5 次三域验证（最多覆盖 25 个 epoch）没有改进时可以提前停止。停止时日志、history 与摘要都会写入 \code{stop_reason=early_stopping}、当前/最佳指标和未提升次数，不再让用户猜测为何没有跑到 100。
此外，\code{max_wall_time_hours=8.5} 使用最近 epoch 耗时预测到下一个安全验证点，只在已验证且
best/last checkpoint 完整写入后停止。它是训练阶段的预测保护，不是操作系统硬截止；数据审计与最终 D4
仍要另留时间，所以报告只给出 4090 整条命令目标约 9 小时的预算估计，不冒充实测。

\section{验证详解}

\subsection{训练和验证的根本差异}

训练需要随机增强、梯度和 optimizer；验证不增强、不建梯度、不开 optimizer。验证数据按 ROI 与训练集
隔离，因此分数衡量的是模型对未见 ROI 的泛化，而不是对训练样本的记忆。

稳定 retrain-v2 的 AutoDL 日志已经提供完整 1292 张 val 证据：EMA 无 TTA 的 JPG SSIM/PSNR 为
0.804480/25.97599，EMA+D4 为 0.809280/26.19677；整条命令耗时 3 小时 25 分，训练部分 2 小时 45 分，
峰值显存 11.63 GiB。v4 尚未在 4090/5090 上完成正式训练，因此这些数字只能作为固定对照，不能写成 v4 已涨分，
更不能伪造平台 test 指标。v4 的实际数字必须由附录中的命令真实生成。

\subsection{SSIM 与 PSNR}

代码用 scikit-image 对每张图单独计算。SSIM 的经典形式\cite{ssim}为

\begin{equation}
  \operatorname{SSIM}(x,y)=
  \frac{(2\mu_x\mu_y+C_1)(2\sigma_{xy}+C_2)}
       {(\mu_x^2+\mu_y^2+C_1)(\sigma_x^2+\sigma_y^2+C_2)}.
\end{equation}

$\mu$ 比较平均亮度，$\sigma$ 比较对比度，$\sigma_{xy}$ 比较结构共同变化。项目数据范围为 1。
PSNR 从均方误差得到：

\begin{equation}
  \operatorname{MSE}=\frac1N\sum_i(\hat y_i-y_i)^2,
  \qquad
  \operatorname{PSNR}=10\log_{10}\frac{1^2}{\operatorname{MSE}}.
\end{equation}

若预测与目标完全相同，MSE 为 0，PSNR 合法地为正无穷；聚合代码会记录无穷值数量。

\subsection{float、uint8 与 JPG 三域}

\begin{description}
  \item[float 域] 直接比较 $[0,1]$ 浮点预测与解码后的浮点目标。它反映模型张量本身的理想质量。
  \item[uint8 域] 预测和目标都按与保存函数相同的 round+clip 规则量化为 $0\sim255$，再除以 255 评分。它衡量 8 位量化损失。
  \item[JPG 域] 只把预测按 ImageSpec 的 RGB 存储模式、quality 100、subsampling 0、optimize false 编码到内存，再解码；目标已经来自官方 JPEG，只做 uint8 量化，不再次压缩。它衡量预测最终落盘可能引入的额外误差。
\end{description}

初赛配置设置 \code{primary_domain=jpg}。这是很合理的本地模型选择代理，因为提交文件就是 JPEG；但在
没有赛事平台评测代码的情况下，不能写成“官方评测器一定逐字逐步完全相同”。正确说法是：JPG 域更
接近最终提交文件形态，官方隐藏 test 分数仍只能由上传 ZIP 后的平台给出。

\subsection{逐图、ROI、最差样本与分层聚合}

每条验证记录包含 target、stem、ROI、organ、坐标、context 可用度、目标图平均值/标准差、activity
proxy 和三域 SSIM/PSNR。聚合层输出：

\begin{itemize}
  \item 图像级 mean/median SSIM、PSNR；
  \item 每个 ROI 内先平均、再在 ROI 间宏平均的 \code{roi_ssim/roi_psnr}；
  \item 最差 10\% 图像的均值，用于发现模型是否只照顾容易样本；
  \item 按 target 的 per-target 指标与跨 target macro；
  \item 按 organ、activity、context availability、图像平均亮度、border/interior 的分层结果；
  \item 本地 proxy：$0.7\times\mathrm{SSIM}+0.3\times$截断归一化 PSNR。
\end{itemize}

当前 manifest 的 organ 是 \code{unknown}，因为目录是 \code{train/marker} 而不是带器官层级的目录；
所以不要把本地结果描述为 colon/liver/stomach 加权评测。坐标解析成功时，Validator 用每个 ROI 实际出现
的最小/最大行列标记 border，其余为 interior。

\subsection{bootstrap 的真实状态}

\term{bootstrap}{自助重采样}通常指反复有放回抽样，以估计指标差的置信区间。项目的
\code{roi_bootstrap_difference()} 会先在每个 ROI 内平均，再对 ROI 做成对重采样，默认 1000 次，输出
候选减基线的均值、95\% 区间和 win/tie/loss。

但是当前 \code{Validator.evaluate()} 虽然读取了 YAML 中
\code{bootstrap_by_roi=true} 的意图，却\textbf{不会在单次普通验证报告中直接生成 bootstrap 区间}。
bootstrap 真正用于 TTA/promotion 比较路径，例如 \code{command_validate()} 的 D4 promotion 和性能 v2
消融门。因而“普通 metrics.json 已经包含 1000 次 bootstrap CI”是不准确的。普通验证包含 ROI 聚合；
候选与基线的置信区间需要走 promotion 比较流程。

\subsection{raw、EMA 与 D4 的选择}

训练期验证对 raw、EMA 分别完整评分，然后按 JPG macro SSIM/PSNR选更优来源。独立
\code{command_validate()} 加载 checkpoint 时，\code{best_jpg} 会读取这个来源；它不是再次自动跑 raw+EMA
两套。若 \code{inference.tta=d4}，命令还会运行无 TTA 与 D4，并利用同一 ROI 清单做 promotion 判断，
把每个 target 的决策写入 \code{validation/tta_decisions.json}。预测时仅在
\code{inference.tta\_policy=validation\_decision} 时按 target 应用该文件；正式 CD68 配置采用
\code{tta\_policy=configured}，因此配置中的 D4 为权威设置，不会被不完整的 promotion 审计覆盖。
本次 seed-2026 最佳 checkpoint 的 JPG 指标由无 TTA 的 0.804480/25.97599 提升至
D4 的 0.809280/26.19677，因而正式提交默认执行 D4。

\section{测试与提交详解}

\subsection{为什么 test 不能显示本地分数}

test 只有 DAPI，没有 CD68 真值。SSIM/PSNR 都属于 full-reference metric，必须同时有预测和真值。
因此运行 \code{autodl-submit} 能得到 1346 张预测和合法 ZIP，却不能得到 test 分数。平台持有隐藏 CD68
标签，上传后才返回 official score。把 val 分数叫“test 分数”或用 val 标签冒充 test 都会误导实验判断。

\subsection{D4、权重来源与 JPEG 保存}

当前推理使用以下四项设置：\path{use_ema=false}、\path{weight_source=best_jpg}、\path{tta=d4}，
以及 \path{tta_policy=configured}。其中 \path{best_jpg} 优先使用 checkpoint 记录的训练期
raw/EMA 胜者；如果旧 checkpoint 没有记录，才根据
\code{use_ema=false} 回退到 raw，且 checkpoint 没有 EMA 时也会安全使用 raw。本次 seed-2026 checkpoint
的训练期胜者是 EMA。完整 1292 张验证集上，D4 把 JPG SSIM 从 0.804480 提高到 0.809280，PSNR 从
25.97599 dB 提高到 26.19677 dB，且 5 个 ROI 的聚合变化均为正，因此正式提交默认执行八次 forward。
这里的“提交 D4”与训练期的“D4 等变一致性”不是同一件事：前者对每张 test 图执行八次前向并平均，
后者在 v4 的 15\% 训练 batch 上追加一次随机变换的一致性约束。提交 D4 不会改变 checkpoint 参数，但训练期
一致性会参与梯度更新。

预测 tensor 先 clamp、乘 255、四舍五入为 uint8，按目标 ImageSpec 保存。当前目标审计为 RGB，所以
最终模式应为 RGB，而不是 L。

\subsection{提交目录与命名}

构建后的真实结构是：

\begin{lstlisting}[style=shell]
outputs/initial_round_max_v4/<run-id>/submission/
|-- results/
|   `-- test/
|       `-- CD68/
|           |-- ROI025_00_00_fake.jpg
|           |-- ...
|           `-- <1346 predictions>
`-- submission_CD68.zip
\end{lstlisting}

提交命名函数会先剥掉 stem 末尾已有的 \code{_fake}，防止产生
\code{_fake_fake.jpg}。ZIP 中每个成员必须从顶层 \code{results/} 开始，不能多包一层任意文件夹。

\subsection{validate-submission 检查什么}

validator 从 test manifest 推导 1346 个期望文件名，与实际目录做集合比较，并逐文件检查：

\begin{itemize}
  \item 缺失、额外和重复输出；
  \item 扩展名和 \code{_fake} 后缀；
  \item 尺寸是否与 manifest 的 $256\times256$ 一致；
  \item JPEG 格式、uint8 dtype、RGB/L mode 和通道数；
  \item 像素范围、NaN/Inf、全黑和全白；近乎常量图给 warning；
  \item ZIP 成员是否都在 \code{results/} 下，是否缺/多文件，CRC 是否通过。
\end{itemize}

校验成功只代表“提交格式正确”，不代表“模型分数高”。格式校验像考试前检查答题卡是否填全；真正得分
还要由平台拿隐藏答案批改。

\section{配置系统：怎样读懂并安全修改 YAML}

\subsection{YAML 是什么}

YAML 是一种供人阅读的层级配置格式。它把“经常调整的参数”从 Python 源码中拿出来，避免每次实验都改
程序。例如：

\begin{lstlisting}[style=yaml]
train:
  epochs: 100
  batch_size: 6
  gradient_accumulation: 2
model:
  context:
    enabled: true
\end{lstlisting}

缩进表示所属关系，不能随意混用 Tab。这里 
\path{train.epochs} 表示最大训练轮数，\path{model.context.enabled=true} 表示请求启用邻域上下文；
v4 还要求 \path{require_verified_grid=true}，审计失败时不会静默猜邻居。
布尔量应写成 \code{true/false}，列表可以写成 \code{[64, 128, 256, 512]}。

\subsection{配置合并优先级}

\code{load_config()} 不是只读一个文件，而是按下列顺序合并；越靠后的值优先级越高：

\begin{enumerate}
  \item \code{configs/default.yaml}：所有实验共享的默认值；
  \item 用户指定的 YAML，例如 \code{configs/initial_round_cd68_max_v4.yaml}；
  \item \code{configs/resolved_local.yaml}：当前文件存在，保存本地审计推导出的数据根、逻辑通道和 worker 设置；
  \item CLI 的 \code{--set key=value}：只覆盖这一次命令，优先级最高。
\end{enumerate}

每个 run 都会把最后生效的配置写入输出目录。复现实验时应保存这个 effective config，而不应只凭记忆记录
“我大概用了哪个 YAML”。例如临时跑 3 epoch 可以写：

\begin{lstlisting}[style=shell]
python -m virtual_staining.cli \
  --log-root log \
  autodl-run \
  --config configs/initial_round_cd68_max_v4.yaml \
  --run-id cd68_max_v4_quickcheck \
  --max-epochs 3
\end{lstlisting}

\important{\code{--max-epochs 3} 是快速检查，不是正式训练。必须换一个新的 run ID，避免把快速实验和
100 epoch 的正式实验混在一起。}

\subsection{当前初赛配置逐项说明}

表~\ref{tab:current-config} 汇总当前真正影响初赛 CD68 训练的主要设置。

\begin{longtable}{p{0.28\textwidth}p{0.18\textwidth}p{0.45\textwidth}}
\caption{\texttt{initial\_round\_cd68\_max\_v4.yaml} 的关键生效配置}\label{tab:current-config}\\
\toprule
配置项 & 当前值 & 含义与影响 \\
\midrule
\endfirsthead
\toprule
配置项 & 当前值 & 含义与影响 \\
\midrule
\endhead
\code{data.root} & AUTO $\rightarrow$ dataset/official & 自动发现官方 train/test；resolved-local 仍可覆盖。\\
\code{data.targets} & CD68 & 当前训练只学习 DAPI 到 CD68。\\
\code{data.num_workers} & 8 & AutoDL 使用 8 个 persistent workers、prefetch 4；Windows 本地可覆盖为 0。\\
logical channels & 1 入 / 1 出 & RGB 文件内容按审计结论压成一个逻辑灰度通道，提交时再恢复 RGB storage。\\
\code{model.name} & \shortstack{multi\_marker\_\\restorer} & 使用四尺度 NAF 主干及 v4 重型三路融合。\\
channels/depths & [64,128,256,512] / [2,2,4,6] & 控制四级 NAF encoder 容量；stem 保持全分辨率。\\
heavy U & [56,88,128,176] & 四尺度完整编解码，局部/轴向大核/空洞卷积，编码深度 [3,4,6,6]。\\
spatial/frequency & 自适应开启 & $1/8$ 根据 patch 学习 FFT low/high 分界与双向调制。\\
multi-route fusion & 四尺度开启 & 高分辨率 NAF/Heavy 门控；$1/8$ NAF/Heavy/frequency softmax + 2 SDPA。\\
context & verified $3\times3$ & tiny shared encoder；$1/4,1/8$ FiLM 与 9-token cross-attention。\\
prototypes & 开启，8 shared + 8 task & 在 $32\times32$ 瓶颈执行余弦注意力与残差融合。\\
task adapters & 开启 & 每个解码尺度使用 CD68 embedding 与零初始化 FiLM/residual adapter。\\
base/detail & direct + zero residual & coarse/detail 零初始化叠加旧 direct logits；calibrator 关闭。\\
loss schedule & two-phase & MSE/Charbonnier/SSIM 等连续调度，深监督退火，末段 float metric proxy。\\
\code{train.epochs} & 100 & 最大轮数；另有 8.5 小时安全停点预测保护。\\
batch / accumulation & 6 / 2 & 有效 batch 仍为 12；给重型激活留出 24 GiB 显存余量。\\
optimizer & grouped fused AdamW & 学习率 $1.8\times10^{-4}$、weight decay $10^{-4}$，bias/norm/residual scale no-decay。\\
scheduler & warmup+cosine & 前 4 epoch 预热，随后余弦衰减。\\
AMP & BF16 & RTX 4090 Tensor Core；SSIM/FFT/prototype 等敏感路径升 float32。\\
EMA & 0.999 & 每个验证点同时比较 raw/EMA，由 JPG 指标选择 checkpoint 来源。\\
validation interval & 5 epoch & 每五轮跑 float、uint8、JPG 三域；最后一轮强制验证。\\
primary domain & JPG & best checkpoint 首先依据更接近提交文件的 JPG 域。\\
inference TTA & D4 configured & 已由稳定对照的全 val/5 ROI 证据确认，提交默认八次变换。\\
\bottomrule
\end{longtable}

\subsection{GPU 加速选项与边界}

当前配置为 CUDA 训练打开了 \code{channels_last}、cuDNN benchmark，并把 float32 matmul precision
设为 \code{high}。固定 $256\times256$ 尺寸适合 cuDNN 自动选择卷积算法。训练器和验证器会把模型与输入
转换为 channels-last 内存格式；当前 inferencer 没有同样的显式转换，因此不能笼统声称“训练、验证、推理
全流程都使用 NHWC”。异步 metric logging、fused AdamW 和 GPU monitor 已开启；\code{torch.compile}
与 profiler 仍关闭，因为前者尚未完成 eager/compiled checkpoint 回归，后者会给正式训练增加诊断开销。
网络中的“并行”表示局部、频率、邻域多路特征在同一计算图中由 GPU 批量 kernel 执行，不等于手写多个
CUDA stream；单卡 autograd 强行多 stream 往往会增加同步和显存。

\subsection{当前通道数的一个重要事实}

官方图片的 JPEG storage mode 是 RGB；审计又证明三通道逐像素相同，所以它在信息意义上是灰度图。
当前 \code{resolved_local.yaml} 明确把输入和四个 target 的逻辑通道设为 1。\code{load_image_tensor()}
读取 RGB 后对三个通道取均值并四舍五入，模型因此构造为 \textbf{1 通道输入、1 通道输出}，参数量与 MAC
统计也以此为准。checkpoint 的 ImageSpec 仍记录原始 RGB storage，保存时把单通道预测复制回三个相同
通道；因此“逻辑单通道”与“磁盘 mode L”不是一回事。

\subsection{新手修改配置的安全原则}

\begin{enumerate}
  \item 每次只改变一个主要因素，并使用新的 run ID；
  \item 先跑 1--3 epoch 检查数据、loss、显存和 checkpoint，再启动最多 100 epoch 的 v4 正式训练；
  \item 比较模型时固定 manifest、seed、训练预算和验证协议；
  \item 不要根据 test 反复调参，因为 test 没有本地标签；
  \item 不要为了“跑起来”把 val 当 test，也不要在 train/val 间混入同一 ROI；
  \item 最终选择必须看 JPG round-trip 指标，而不能只看训练 loss。
\end{enumerate}

\section{项目文件结构与模块职责}

\subsection{从根目录看项目}

\begin{lstlisting}[style=shell]
project/
|-- configs/                  # YAML configurations
|-- dataset/                  # data only; never import it into source code
|   `-- official/{train,test}
|-- src/virtual_staining/     # installable Python package
|-- docs/essay/report.tex     # the single editable technical guide/report
|-- artifacts/                # manifests and data audit artifacts
|-- outputs/                  # checkpoints, validation and predictions
|-- output/pdf/               # compiled technical report PDF
|-- log/                      # compact AutoDL logs and downloadable bundles
|-- pyproject.toml            # package metadata and Python dependencies
`-- README.md                 # project entry point for users
\end{lstlisting}

\important{按上传精简要求，当前工作区没有 \code{tests/} 目录，\code{.ruff\_cache/} 也已清理；上面的树
因此不列测试资产。本报告不声称该精简副本通过了 pytest。若此版本继续长期开发，应在独立开发分支恢复
与当前接口匹配的回归测试，而不是把训练数据、checkpoint 或提交结果当成测试替代品。}

源代码采用 \term{src-layout}{包放在 \code{src/} 下的 Python 工程布局}。运行
\code{pip install -e .} 后，Python 能从任意工作目录找到 \code{virtual_staining}；\code{-e} 表示 editable，
修改源码后不必重新复制安装。若完全不安装项目，也可以在项目根目录设置
\code{PYTHONPATH=src}，但对新手更容易漏设，因此推荐 editable install。

\subsection{包入口与公共模块}

\begin{longtable}{p{0.40\textwidth}p{0.52\textwidth}}
\caption{\texttt{src/virtual\_staining} 顶层模块}\label{tab:top-modules}\\
\toprule
文件 & 主要职责 \\
\midrule
\endfirsthead
\toprule 文件 & 主要职责 \\
\midrule
\endhead
\code{__init__.py} & 定义包版本和公开入口，使目录成为 Python package。\\
\code{constants.py} & 集中维护 DAPI、四个目标 marker、别名与文件扩展名等稳定常量。\\
\code{config.py} & 读取/深度合并 YAML，解析 CLI 覆盖，严格检查字段和取值，输出 effective config。\\
\code{cli.py} & 全部命令的总控制器；把数据发现、训练、验证、预测、提交和 AutoDL 日志串成可执行工作流。\\
\bottomrule
\end{longtable}

\subsection{数据模块}

\begin{longtable}{p{0.38\textwidth}p{0.54\textwidth}}
\caption{\texttt{src/virtual\_staining/data} 模块}\label{tab:data-modules}\\
\toprule 文件 & 主要职责 \\
\midrule
\endfirsthead
\toprule 文件 & 主要职责 \\
\midrule
\endhead
\code{activity.py} & 计算/分箱图像活性，用于分层报告或 activity-aware sampler。\\
\code{audit.py} & 全量检查尺寸、mode、通道差、值域、配对、相关性和对齐，生成数据审计报告。\\
\code{dataset.py} & 根据 manifest 加载 DAPI/目标，转换为 CHW float tensor，支持单/多目标及邻域字段。\\
\code{discovery.py} & 自动搜索候选数据根，识别 split/organ/marker 层级并给候选打分。\\
\code{grouped_folds.py} & 以真实 ROI/group 为单位构建交叉验证 fold，防止 patch 级随机切分泄漏。\\
\code{loader_helpers.py} & 安全构造 DataLoader、worker seed、pin-memory/prefetch 等参数。\\
\code{manifest.py} & 配对 DAPI 与各 marker、解析 canonical key/hash/group，生成 train/val/test CSV 和泄漏报告。\\
\code{neighborhood.py} & 用行列坐标查找 $3\times3$ 邻居，返回 tiles、有效 mask 与相对 offset。\\
\code{roi_audit.py} & 检查二维网格方向、空洞、重复坐标、边界连续性以及跨 split 邻域风险。\\
\code{roi_index.py} & 严格解析 \code{ROI000_00_00} 为 ROI、row、col，并建立坐标索引。\\
\code{samplers.py} & 实现均匀或活性相关采样策略。\\
\code{transforms.py} & 对 DAPI 和标签同步执行几何增强，对 DAPI 单独执行允许的强度增强。\\
\code{__init__.py} & 汇总数据子包的公开符号。\\
\bottomrule
\end{longtable}

\subsection{模型模块}

\begin{longtable}{p{0.38\textwidth}p{0.54\textwidth}}
\caption{\texttt{src/virtual\_staining/models} 模块}\label{tab:model-modules}\\
\toprule 文件 & 主要职责 \\
\midrule
\endfirsthead
\toprule 文件 & 主要职责 \\
\midrule
\endhead
\code{baseline_unet.py} & 轻量 Residual U-Net 基线，用于 smoke 和回归验证。\\
\code{camp_vs_v2.py} & 可选 v2 总模型，连接 local encoder、global mixer、conditioning、decoder、base/detail 和 calibrator；当前初赛 YAML 不选择它。\\
\code{naf_blocks.py} & LayerNorm2d、SimpleGate、NAF block 等卷积恢复基础积木。\\
\code{naf_local_encoder.py} & 四级局部编码器，输出 $1/2$ 至 $1/16$ 多尺度特征。\\
\code{restoration_transformer.py} & GPU 友好的 Restormer-lite 通道注意力和门控前馈模块。\\
\code{task_organ_conditioning.py} & marker/organ embedding 与零初始化 FiLM 条件调制。\\
\code{context_tile_encoder.py} & 共享参数编码邻居 tile，结合位置 offset 和有效 mask。\\
\code{context_fusion.py} & 将邻域上下文以残差 FiLM 或瓶颈交叉注意力注入中心特征。\\
\code{lightweight_detail_unet.py} & 四尺度大核 depthwise/gated 轻量 U 分支，向 $1,1/2,1/4$ skip 注入零初始化细节残差。\\
\code{heavy_detail_unet.py} & v4 四尺度重型 U 分支；高分辨率使用 depthwise/local gated block，低分辨率组合局部、轴向大核与空洞卷积。\\
\code{spatial_frequency_mixer.py} & 在 $1/8$ 瓶颈执行 float32 FFT low/high 分解；v4 由当前图像预测受限 cutoff，让低/高频相互调制，并以 route-only 接口避免第三路重复携带 NAF identity。\\
\code{multi_route_fusion.py} & 对 NAF、重型 U 与频率路线同时计算逐通道/逐空间 softmax gate，并以小残差系数和两层低分辨率 SDPA 注入。\\
\code{hierarchical_prototypes.py} & 多尺度 shared/marker/organ 原型、用量与 dead-prototype 诊断。\\
\code{prototype_mixer.py} & 旧版余弦原型匹配核心，供兼容模型和扩展模块复用。\\
\code{laplacian_decoder.py} & 预测 coarse/detail，并在 v3/v4 中以零初始化 residual 叠加已验证的 direct logits，避免初始输出突变。\\
\code{intensity_calibrator.py} & 学习受限的 per-target gain/bias，校正整体强度而避免无界漂移。\\
\code{multi_marker_restorer.py} & \textbf{当前有效主模型}：四尺度 NAF 编解码、重型 Detail U、自适应频率路线、多尺度三路融合、verified context、shared/task prototypes、task adapters、稳定 base/detail 与三尺度输出。\\
\code{dapi_mae.py} & 可选 DAPI masked autoencoder 预训练，不使用标签并需 fold-local 防泄漏。\\
\code{registry.py} & 按配置名构造 Residual U-Net、MultiMarkerRestorer 或 CAMP-VS v2，并统一输出接口。\\
\code{__init__.py} & 汇总模型子包接口。\\
\bottomrule
\end{longtable}

\subsection{损失与指标模块}

\begin{longtable}{p{0.38\textwidth}p{0.54\textwidth}}
\caption{损失和指标模块}\label{tab:loss-metric-modules}\\
\toprule 文件 & 主要职责 \\
\midrule
\endfirsthead
\toprule 文件 & 主要职责 \\
\midrule
\endhead
\code{losses/charbonnier.py} & 平滑且较抗离群值的像素重建损失。\\
\code{losses/ssim.py} & 可微 SSIM/MS-SSIM；敏感运算强制 float32。\\
\code{losses/gradient.py} & 用 Sobel 梯度约束边缘结构。\\
\code{losses/frequency.py} & 在 log-frequency 域比较纹理能量。\\
\code{losses/pyramid.py} & 多尺度 Laplacian/pyramid 重建约束。\\
\code{losses/statistics.py} & 约束均值、标准差等强度统计。\\
\code{losses/competition_proxy.py} & 后段训练使用的可微 float SSIM/PSNR 代理；不冒充 uint8/JPEG round-trip。\\
\code{losses/prototype.py} & prototype activation/diversity 等辅助约束。\\
\code{losses/shift_tolerant.py} & 可选小位移容忍损失；当前对齐审计为 0，默认关闭。\\
\code{losses/composite.py} & v1 固定权重复合损失和深监督汇总。\\
\path{losses/scheduled_composite.py} & v2 two-phase 权重插值与所有损失项统一入口。\\
\code{metrics/image_metrics.py} & 逐图计算 scikit-image SSIM/PSNR。\\
\code{metrics/domain_metrics.py} & float、uint8、JPG 三域量化与 ROI bootstrap 差值。\\
\code{metrics/aggregation.py} & 汇总 mean/median、ROI、target、macro、worst-10\% 和 proxy。\\
\code{metrics/promotion.py} & 按置信区间、分层退化和泄漏状态判断候选能否晋级。\\
各目录 \code{__init__.py} & 维护稳定导入接口。\\
\bottomrule
\end{longtable}

\subsection{训练、推理与实验引擎}

\begin{longtable}{p{0.34\textwidth}p{0.58\textwidth}}
\caption{\texttt{src/virtual\_staining/engine} 模块}\label{tab:engine-modules}\\
\toprule 文件 & 主要职责 \\
\midrule
\endfirsthead
\toprule 文件 & 主要职责 \\
\midrule
\endhead
\code{common.py} & 公用 DataLoader、模型/损失/优化器构造和配置辅助逻辑。\\
\code{optim.py} & 构建 no-decay 参数组与 fused AdamW，并记录实际采用的优化器实现。\\
\code{trainer.py} & 训练主循环：BF16 AMP、D4 一致性、EMA、OOM microbatch、GPU 监控、时间保护、验证、early stopping 和保存。\\
\code{validator.py} & 三域逐图推理、分组聚合、raw/EMA 比较和验证产物写入。\\
\code{inferencer.py} & 无标签确定性预测、D4 逆变换、显存降级和 JPEG 输出。\\
\code{checkpoint.py} & 保存/恢复模型、EMA、optimizer、scheduler、scaler、RNG、manifest/config hash。\\
\code{ema.py} & 参数指数滑动平均的更新、state dict 与临时应用。\\
\code{ensemble.py} & 多 checkpoint 预测组合与通用集成入口。\\
\code{ensemble_optimizer.py} & 仅用 OOF/val 预测寻找非负权重，失败时可回退均匀平均。\\
\code{model_soup.py} & 对兼容、同血缘模型做权重 soup，并检查安全前提。\\
\code{multistage.py} & 管理 multitask、target、organ、metric 等多阶段训练与恢复。\\
\code{multitask_optimizer.py} & equal/FAMO 类多任务权重及梯度冲突处理实验。\\
\code{pretrainer.py} & DAPI-MAE 的 fold-local 预训练流程。\\
\code{experiment_registry.py} & 记录实验父子关系、配置、状态和指标，防止选择性汇报。\\
\code{ablation_budget.py} & 定义 smoke/screen/confirm/full 预算与 A0--A3 串行约束。\\
\code{complexity_report.py} & 统计参数、近似 MAC 和模块占比。\\
\code{gradient_monitor.py} & 监控不同任务梯度范数/夹角，诊断负迁移。\\
\code{prototype_diagnostics.py} & 导出 prototype 使用率、熵、相似度等诊断。\\
\code{prototype_monitor.py} & 训练期跟踪 dead/collapsed prototype，可选生成注意力图。\\
\code{__init__.py} & 汇总引擎子包接口。\\
\bottomrule
\end{longtable}

\subsection{提交与工具模块}

\begin{longtable}{p{0.34\textwidth}p{0.58\textwidth}}
\caption{提交和工具模块}\label{tab:utility-modules}\\
\toprule 文件 & 主要职责 \\
\midrule
\endfirsthead
\toprule 文件 & 主要职责 \\
\midrule
\endhead
\code{submission/writer.py} & 按竞赛目录/命名复制预测，避免二次 JPEG 编码，并创建 ZIP。\\
\code{submission/validator.py} & 校验数量、命名、mode、尺寸、像素异常、ZIP 路径和 CRC。\\
\code{submission/__init__.py} & 暴露 submission 公共接口。\\
\code{utils/device.py} & 发现 CPU/CUDA、显存与 BF16 能力，解析设备和硬件 profile。\\
\code{utils/gpu_monitor.py} & 周期采样 GPU 利用率、显存、功耗与温度，按 epoch 汇总 mean/p95/max。\\
\code{utils/image_io.py} & 统一 RGB/L 读取、float/uint8 转换和高质量 JPEG 保存。\\
\code{utils/logging.py} & 创建 console/file logger、结构化 JSON 日志和运行元数据。\\
\code{utils/paths.py} & 使用 pathlib 处理中文、空格、相对和绝对路径。\\
\code{utils/seed.py} & 固定 Python、NumPy、Torch 和 DataLoader 随机种子。\\
\code{utils/__init__.py} & 汇总工具子包接口。\\
\bottomrule
\end{longtable}

\subsection{outputs、artifacts 与 log 的区别}

\begin{itemize}
  \item \code{artifacts/} 主要是“数据事实”：发现结果、manifest、审计、泄漏报告和校验和；
  \item \code{outputs/} 主要是“模型产物”：checkpoint、验证逐图 CSV/JSON、预测和 submission；
  \item \code{log/} 主要是“方便人和下一位工程师阅读的轻量记录”：终端日志、环境、GPU、吞吐、摘要及可下载 ZIP。
\end{itemize}

不要把三者混为一谈。删除 checkpoint 可以释放空间，但 manifest 与审计文件很小且决定实验可复现性，通常
应该保留。上传给协作者做性能诊断时，优先提供 log bundle、effective config 和 metrics，而不是上传数 GB
的所有 checkpoint；只有需要复现权重行为时才提供选中的 best checkpoint。

\section{总结：从一张 DAPI 到竞赛 ZIP}

CAMP-VS 把虚拟染色组织成一条可审计而不是“只会跑一个网络”的流水线：先发现数据并把同名 DAPI/CD68
配成 manifest；严格按 ROI 隔离 train/val；把 RGB storage 转换成单逻辑通道 float tensor；当前 v4 用
四尺度 NAF 主干提取稳定结构，同时由重型 Heavy Omni U-Net 恢复多尺度细胞细节、由自适应 FFT 独立路线
分解并双向调制 low/high 频率，真实 $3\times3$ DAPI 邻域提供 ROI 上下文。四个尺度使用逐通道与逐空间
softmax gate 融合，$1/8$ 瓶颈再用两层原生 SDPA、context 和 shared/task cosine prototypes 建立全局语义；
解码端以 CD68 adapter 和 zero-init base/detail residual 恢复三尺度预测；最后在 float、uint8、JPG 三个域
验证 raw/EMA，并以配置化 D4 推理官方 test，进行严格 ZIP 校验。

项目当前有三个特别重要的边界。第一，本地 val 有 CD68 真值，所以能够计算 SSIM/PSNR；官方 test 只有
DAPI，本地不能算分。第二，文件存储为 RGB，但当前模型实际使用 1 输入/1 输出逻辑通道，提交时才恢复
RGB storage。第三，v4 的重型 U、自适应 spatial/frequency、多尺度三路融合、两层低分辨率 SDPA、verified
context、cross-attention、prototype、task adapter 和 residual base/detail 确实进入 forward；完整
Restormer/Mamba、diffusion、GAN、calibrator、FAMO 与 shift-tolerant 等替代模块仍关闭。报告不会把“代码存在”
写成“默认已使用”。

对初学者而言，最重要的实验习惯仍是保证比较公平：固定 split、seed、训练预算、三域协议和 JPEG 参数；
保存 effective config、GPU/模块日志和逐图指标。v4 是“单次重训的候选组合”，不是已证明上涨的结果。
只有它的 D4 JPG 指标稳定超过固定对照 0.809280/26.19677、各 ROI/活动度/边界分层没有集中退化，且平台
分数也提高后，才能声称性能升级成功。

\appendix

\section{AutoDL 从上传到训练的保姆级流程}

本附录假设项目已上传到 \code{/root/autodl-tmp/project}，官方数据位于项目内
\code{dataset/official/train} 和 \code{dataset/official/test}，环境名为 \code{MEDICAL}。如果实际项目路径
不同，只改第一条 \code{cd}；不要修改 Python 源码里的路径。

\subsection{第 0 步：确认上传目录}

\begin{lstlisting}[style=shell]
cd /root/autodl-tmp/project
pwd
ls
ls dataset/official/train/DAPI | head
ls dataset/official/test/DAPI | head
\end{lstlisting}

根目录至少应看到 \code{pyproject.toml}、\code{src}、\code{configs}、\code{dataset}。train 下应有 DAPI
和 CD68 等 marker；test 下只有 DAPI 是正常的。

\subsection{第 1 步：激活环境并让 Python 找到项目}

推荐 editable install：

\begin{lstlisting}[style=shell]
conda activate MEDICAL
python -m pip install -e . --no-deps
\end{lstlisting}

\code{-e} 不会复制一份大型项目，只在环境中登记源码位置；\code{--no-deps} 表示依赖已经装好，不再下载
PyTorch 等大包。若确实不愿登记项目，每次新终端都必须执行：

\begin{lstlisting}[style=shell]
conda activate MEDICAL
export PYTHONPATH="$PWD/src:$PYTHONPATH"
\end{lstlisting}

然后验证导入和 GPU：

\begin{lstlisting}[style=shell]
unset OMP_NUM_THREADS
export OMP_NUM_THREADS=1
export MKL_NUM_THREADS=1
python -c "import virtual_staining; print('package ok')"
python -c "import torch; print(torch.__version__, torch.version.cuda); \
print('cuda=', torch.cuda.is_available()); \
print(torch.cuda.get_device_name(0) if torch.cuda.is_available() else 'CPU')"
\end{lstlisting}

若 \code{cuda=False}，不要启动 100 epoch。先解决 PyTorch/CUDA 环境，否则模型会落到 CPU，训练可能从
数小时变成数天。

\subsection{第 2 步：可选的 3 epoch 流程检查}

第一次上传建议先用独立 run ID 做快速检查：

\begin{lstlisting}[style=shell]
python -m virtual_staining.cli --log-root log autodl-run \
  --config configs/initial_round_cd68_max_v4.yaml \
  --data-root AUTO \
  --target CD68 \
  --run-id cd68_max_v4_quickcheck \
  --max-epochs 3
\end{lstlisting}

终端会显示 tqdm 进度条和 epoch/batch 进度。这个命令会真实读取、审计、训练、保存和验证，但 3 epoch
结果只证明流程打通。它既不是正式成绩，也不能取代完整训练。

\subsection{第 3 步：v4 正式训练与验证}

使用一个新的、永不复用的 run ID：

\begin{lstlisting}[style=shell]
python -m virtual_staining.cli --log-root log autodl-run \
  --config configs/initial_round_cd68_max_v4.yaml \
  --data-root AUTO \
  --target CD68 \
  --run-id cd68_max_v4_seed2026
\end{lstlisting}

一条命令依次执行环境记录、数据发现、manifest、ROI 邻域审计、最多 100 epoch 训练、raw/EMA 三域验证和
最终 D4 验证。v4 每 5 轮验证，15\% batch 执行训练期等变分支，训练阶段有 8.5 小时预测式安全停点。
基于 v3 远程吞吐与 v4 75.997 GMAC 的工程估算是后段约 4.5--5.2 分钟/轮，整条命令目标 9 小时左右，但尚未在 AutoDL 完整实测。
不要关闭
SSH 页面后误以为进程一定继续；若需要断线运行，应在 AutoDL 的 tmux 中执行，或使用平台提供的后台任务。

实时观察 GPU 和日志可另开一个终端：

\begin{lstlisting}[style=shell]
watch -n 2 nvidia-smi
find log/cd68_max_v4_seed2026 -maxdepth 3 -type f
\end{lstlisting}

具体日志文件名以该 run 的 \code{log/} 目录实际内容为准；\code{find log -maxdepth 3 -type f} 可列出全部。
如果 GPU 利用率偶尔下降但进度条仍前进，可能正在做 JPEG 验证或加载数据，并不等于卡死。长时间为 0\%
且 batch 不增长时，再检查日志末尾和系统内存。

\subsection{第 4 步：读取验证分数}

完整训练结束后，先查看：

\begin{lstlisting}[style=shell]
python -c "import json; \
m=json.load(open('outputs/initial_round_max_v4/cd68_max_v4_seed2026/validation/metrics.json')); \
j=m['tta_comparison']['d4']['domains']['jpg']['macro']; \
print('D4 JPG SSIM =', j['mean_ssim']); \
print('D4 JPG PSNR =', j['mean_psnr']); \
print('fixed baseline = 0.809280 / 26.19677')"
\end{lstlisting}

这一步使用 official train 中隔离出的 val：输入和 CD68 真值都存在，因此可以算分。以后改模型时，用相同
manifest、seed、训练预算和这条读取方法比较；新版本 JPG SSIM/PSNR 变高是“本地证据”，还需要逐图和
ROI 分层确认没有集中退化。它并不等于排行榜一定上升。

如需对某个 checkpoint 单独重跑验证，可执行：

\begin{lstlisting}[style=shell]
python -m virtual_staining.cli --log-root log validate \
  --config configs/initial_round_cd68_max_v4.yaml \
  --data-root AUTO \
  --checkpoint outputs/initial_round_max_v4/cd68_max_v4_seed2026/checkpoints/best_ssim.ckpt
\end{lstlisting}

\subsection{第 5 步：对 test 推理并生成提交包}

最稳妥的做法是显式写出刚才选中的 checkpoint：

\begin{lstlisting}[style=shell]
python -m virtual_staining.cli --log-root log autodl-submit \
  --config configs/initial_round_cd68_max_v4.yaml \
  --data-root AUTO \
  --target CD68 \
  --checkpoint outputs/initial_round_max_v4/cd68_max_v4_seed2026/checkpoints/best_ssim.ckpt
\end{lstlisting}

也可以不写 checkpoint，让程序在 \code{outputs/**/best_ssim.ckpt} 中按修改时间自动选最新文件：

\begin{lstlisting}[style=shell]
python -m virtual_staining.cli --log-root log autodl-submit \
  --config configs/initial_round_cd68_max_v4.yaml \
  --data-root AUTO \
  --target CD68
\end{lstlisting}

当前 \code{build_parser()} 没有给 \code{autodl-submit} 注册 \code{--run-id} 参数，所以不要照旧版说明在
提交命令末尾添加它，否则 argparse 会直接报告“unrecognized arguments”。自动选择在有多个实验时可能拿到
错误权重，因此正式提交推荐第一种显式 checkpoint 写法。提交输出目录由 checkpoint 的父 run 目录决定。

test 没有 CD68 真值，因此命令只生成预测并检查格式，不会显示真实 SSIM/PSNR。最终上传：

\begin{lstlisting}[style=shell]
outputs/initial_round_max_v4/cd68_max_v4_seed2026/submission/submission_CD68.zip
\end{lstlisting}

竞赛平台用隐藏标签评分后，返回的才是 official test score。

\subsection{第 6 步：下载哪些文件}

\begin{longtable}{p{0.49\textwidth}p{0.43\textwidth}}
\caption{建议从 AutoDL 下载的产物}\label{tab:download-files}\\
\toprule 文件 & 用途 \\
\midrule
\endfirsthead
\toprule 文件 & 用途 \\
\midrule
\endhead
\code{log/<run>/...log_bundle.zip} & 首选诊断包；包含环境、配置、阶段耗时、GPU/资源、指标和日志摘要。\\
\code{outputs/.../validation/metrics.json} & 三域汇总指标。\\
\code{outputs/.../validation/per_image.csv} & 找最差样本、ROI 和亮度分层退化。\\
\code{outputs/.../pipeline_report.json} & 一次流水线的完成状态和关键路径。\\
\path{outputs/.../checkpoints/best_ssim.ckpt} & 继续训练或重新推理需要；文件较大，不做诊断时可不传给协作者。\\
\path{outputs/.../submission/submission_CD68.zip} & 上传竞赛平台的最终包。\\
\bottomrule
\end{longtable}

如果要交给另一位工程师阅读并改进，至少提供日志压缩包、汇总/逐图指标、effective config 和
平台分数截图/文本。没有这些 provenance，只给一句“跑了 15 小时分数不高”，很难区分瓶颈来自数据加载、
模型、loss、验证编码还是环境。

\section{常见问题与排错顺序}

\begin{longtable}{p{0.28\textwidth}p{0.64\textwidth}}
\caption{新手常见故障}\label{tab:troubleshooting}\\
\toprule 现象 & 排查与处理 \\
\midrule
\endfirsthead
\toprule 现象 & 排查与处理 \\
\midrule
\endhead
\code{No module named virtual_staining} & 确认位于项目根目录，执行 \code{pip install -e . --no-deps}，或正确设置 \code{PYTHONPATH="$PWD/src:$PYTHONPATH"}。\\
CUDA availability 为 False & 检查租用实例是否有 GPU、驱动、PyTorch CUDA wheel；修好前不要 CPU 硬跑正式训练。\\
CUDA out of memory & 训练器先自动把 v4 物理 batch 6 切成更小 chunk，effective batch 不变；仍失败时用 \code{--set train.batch_size=3 --set train.gradient_accumulation=4} 保持有效 batch 12，再不够才开启 heavy-U activation checkpoint。\\
进度条似乎不动 & 看 \code{nvidia-smi}、系统内存、日志末尾；数据审计不是训练 batch。验证虽没有逐图 tqdm，但会输出 raw/EMA 或 D4 pass 的开始/完成、样本数和耗时；若连 pass 开始信息也没有，再排查进程。\\
DataLoader worker 崩溃 & Linux 先减少 \code{data.num_workers}；Windows 本地调试设为 0。检查共享内存、损坏图和 OOM，而非盲目重启。\\
JPG 指标低于 float & 量化和 JPEG 会损失细节，这是三域评估的目的；确认质量 100、subsampling 0 且没有二次编码。\\
val 很好、平台分数差 & 检查 ROI 分布差异、过拟合、submission target/命名、raw/EMA/TTA 选择以及本地指标与官方公式差异。\\
run ID 已存在 & 使用新的唯一 run ID；不要覆盖旧 run 后再拿混合文件做比较。\\
磁盘迅速增大 & 保留 best checkpoint、日志和指标，清理失败/重复 run 的大 checkpoint 前先核对路径；不要删除官方数据和唯一 best 权重。\\
\bottomrule
\end{longtable}

\section{实现事实核对清单}

表~\ref{tab:fact-check} 把容易混淆的说法与当前代码真实状态并列，便于以后代码升级后重新审阅报告。

\begin{longtable}{p{0.40\textwidth}p{0.52\textwidth}}
\caption{截至本报告生成时的实现事实}\label{tab:fact-check}\\
\toprule 常见说法 & 当前真实状态 \\
\midrule
\endfirsthead
\toprule 常见说法 & 当前真实状态 \\
\midrule
\endhead
“训练集 6296 张” & 6296 是每个 marker 的总配对数；当前 manifest 为 5004 train、1292 val。\\
“测试集可计算 SSIM” & 错。1346 张 test 只有 DAPI，真实标签在平台。\\
“图像是 L 单通道” & storage 仍为 RGB；resolved-local 让模型实际使用 1 入/1 出逻辑通道，提交时恢复 RGB。\\
“v4 用了 $3\times3$ 邻域” & 对。只接受 filename 坐标和 verified grid；train/val ROI 与邻接边均不重叠。\\
“默认用了 prototype” & 对。当前启用 8 个 shared 与 8 个 CD68 task prototypes，并有两项 0.001 正则。\\
“NAF 模型完全没有任何激活” & NAF block 的核心是 activation-free；全模型其他分支仍可包含 GELU 等操作。\\
“EMA 一定比 raw 好” & 错。每个验证点同时评分，由 JPG 指标选择；短训练 EMA 可能滞后。\\
“验证只按 YAML domains 列表算一域” & 当前 Validator 实际总是生成 float、uint8、JPG 三域。\\
“JPG 域等于官方评分” & 它最接近提交文件，但官方公式/隐藏处理不在本地代码中，不能宣称完全一致。\\
“普通 metrics.json 含 bootstrap CI” & 普通验证有 ROI 聚合；成对 CI 在 promotion/TTA 比较路径生成。\\
“v4 一定跑满 100 epoch” & 错。每 5 epoch 验证；5 个验证事件无改进或预测将超过 8.5 小时安全停点都可提前停止。\\
“所有阶段都用了 channels-last” & Trainer/Validator 显式使用；当前 Inferencer 没有同样的显式转换。\\
“v4 是一个完整 Restormer 或 Mamba” & 错。它是 29,130,894 参数、约 75.997 GMAC 的 NAF、重型 U、自适应 FFT 与 context 三路模型；只有 $1/8$ 瓶颈的融合器使用两层原生 SDPA，没有完整 Restormer/Mamba forward。\\
“base/detail 会随机替换旧输出头” & v3/v4 都不会。direct logits 保留，coarse/detail 只作为零初始化 residual，初始 prediction 与 direct head 逐位一致。\\
“训练 loss 里的 metric proxy 等于 JPG” & 错。它只对齐 float SSIM/PSNR；uint8/JPEG round-trip 只在验证阶段真实执行。\\
“OOM 后 accumulation 会自动翻倍” & 错。v4 只切分同一 DataLoader batch，并在必要时重启累积窗口，保持配置 effective batch。\\
\bottomrule
\end{longtable}

\begin{thebibliography}{99}

\bibitem{unet}
O. Ronneberger, P. Fischer, and T. Brox,
``U-Net: Convolutional Networks for Biomedical Image Segmentation,''
\emph{Medical Image Computing and Computer-Assisted Intervention}, 2015.
\url{https://arxiv.org/abs/1505.04597}

\bibitem{nafnet}
L. Chen, X. Chu, X. Zhang, and J. Sun,
``Simple Baselines for Image Restoration,''
\emph{European Conference on Computer Vision}, 2022.
\url{https://arxiv.org/abs/2204.04676}

\bibitem{restormer}
S. W. Zamir et al.,
``Restormer: Efficient Transformer for High-Resolution Image Restoration,''
\emph{IEEE/CVF Conference on Computer Vision and Pattern Recognition}, 2022.
\url{https://openaccess.thecvf.com/content/CVPR2022/html/Zamir_Restormer_Efficient_Transformer_for_High-Resolution_Image_Restoration_CVPR_2022_paper.html}

\bibitem{ssim}
Z. Wang, A. C. Bovik, H. R. Sheikh, and E. P. Simoncelli,
``Image Quality Assessment: From Error Visibility to Structural Similarity,''
\emph{IEEE Transactions on Image Processing}, vol. 13, no. 4, pp. 600--612, 2004.
\url{https://doi.org/10.1109/TIP.2003.819861}

\bibitem{adamw}
I. Loshchilov and F. Hutter,
``Decoupled Weight Decay Regularization,''
\emph{International Conference on Learning Representations}, 2019.
\url{https://arxiv.org/abs/1711.05101}

\bibitem{pytorch_adamw}
PyTorch Contributors,
``AdamW---PyTorch Documentation.''
\url{https://docs.pytorch.org/docs/stable/generated/torch.optim.AdamW}

\bibitem{pytorch_amp}
PyTorch Contributors,
``Automatic Mixed Precision Package---PyTorch Documentation.''
\url{https://docs.pytorch.org/docs/stable/accelerator/amp.html}

\bibitem{thermofisher_dapi}
Thermo Fisher Scientific,
``DAPI Stain.''
\url{https://www.thermofisher.com/us/en/home/life-science/cell-analysis/fluorophores/dapi-stain.html}

\bibitem{hpa_cd68}
The Human Protein Atlas,
``CD68 tissue expression and primary data.''
\url{https://www.proteinatlas.org/ENSG00000129226-CD68/tissue/primary+data}

\bibitem{cd68_specificity}
R. Gottfried et al.,
``Expression of CD68 in non-myeloid cell types,''
\emph{Scandinavian Journal of Immunology}, vol. 67, no. 5, pp. 453--463, 2008.
\url{https://pubmed.ncbi.nlm.nih.gov/18405323/}

\bibitem{memory_t_cells}
D. L. Farber, N. A. Yudanin, and N. P. Restifo,
``Human memory T cells: generation, compartmentalization and homeostasis,''
\emph{Nature Reviews Immunology}, vol. 14, pp. 24--35, 2014.
\url{https://pmc.ncbi.nlm.nih.gov/articles/PMC4032067/}

\bibitem{hla_dr_review}
L. J. Stern and J. M. Calvo-Calle,
``HLA-DR: molecular insights and vaccine design,''
\emph{Current Pharmaceutical Design}, vol. 15, no. 28, pp. 3249--3261, 2009.
\url{https://pmc.ncbi.nlm.nih.gov/articles/PMC3615543/}

\bibitem{hpa_vimentin}
The Human Protein Atlas,
``VIM protein expression summary.''
\url{https://www.proteinatlas.org/ENSG00000026025-VIM}

\bibitem{fluorescence_quantitation}
J. C. Waters,
``Accuracy and precision in quantitative fluorescence microscopy,''
\emph{Journal of Cell Biology}, vol. 185, no. 7, pp. 1135--1148, 2009.
\url{https://pmc.ncbi.nlm.nih.gov/articles/PMC2712964/}

\bibitem{convir}
Y. Cui et al.,
``Revitalizing Convolutional Network for Image Restoration,''
\emph{IEEE Transactions on Pattern Analysis and Machine Intelligence}, 2024.
\url{https://github.com/c-yn/ConvIR}

\bibitem{darkir}
G. Feijoo et al.,
``DarkIR: Robust Low-Light Image Restoration,''
\emph{IEEE/CVF Conference on Computer Vision and Pattern Recognition}, 2025.
\url{https://github.com/cidautai/DarkIR}

\bibitem{esc}
D. Lee et al.,
``ESC: Emulating Self-Attention with Convolution for Efficient Image Super-Resolution,''
\emph{IEEE/CVF International Conference on Computer Vision}, 2025.
\url{https://github.com/dslisleedh/ESC}

\bibitem{oknet}
Y. Cui et al.,
``Omni-Kernel Network for Image Restoration,''
\emph{AAAI Conference on Artificial Intelligence}, 2024.
\url{https://github.com/c-yn/OKNet}

\bibitem{adair}
Y. Cui et al.,
``AdaIR: Adaptive All-in-One Image Restoration via Frequency Mining and Modulation,''
\emph{International Conference on Learning Representations}, 2025.
\url{https://github.com/c-yn/AdaIR}

\bibitem{eqreg}
Y. Bai et al.,
``A Regularization-Guided Equivariant Approach for Image Restoration,''
\emph{IEEE/CVF Conference on Computer Vision and Pattern Recognition}, 2025.
\href{https://openaccess.thecvf.com/content/CVPR2025/html/Bai_A_Regularization-Guided_Equivariant_Approach_for_Image_Restoration_CVPR_2025_paper.html}{CVPR 2025 open-access paper}

\end{thebibliography}

\end{document}
